\documentclass[11pt,a4paper]{article}

% Fonts and language (pdfLaTeX)
\usepackage[utf8]{inputenc}
\usepackage[T1, T2A]{fontenc}
\usepackage[english, russian]{babel}

% Page geometry
\usepackage[margin=2.5cm]{geometry}

% Math
\usepackage{amsmath, amssymb, amsthm, mathtools}

% Tables
\usepackage{booktabs}
\usepackage{array}

% Graphics
\usepackage{graphicx}
\graphicspath{{figures/}}
\usepackage{float}

% Hyperlinks
\usepackage{xcolor}
\usepackage[colorlinks=true,
            linkcolor=blue!50!black,
            citecolor=blue!50!black,
            urlcolor=blue!50!black]{hyperref}

% Bibliography
\usepackage[numbers, sort&compress]{natbib}
\bibliographystyle{plainnat}

% Typography
% \usepackage{microtype}

% Title
\title{Арифметическая структура стандартного генетического кода}
\author{Руслан Хафизов}
\date{\today\\[1ex]
\small DOI препринта:~\href{https://doi.org/10.5281/zenodo.20751290}{10.5281/zenodo.20751290}}
% Переносимые имена файлов: разрыв строки допускается после подчёркивания
\let\origunderscore\_
\renewcommand{\_}{\origunderscore\allowbreak}
\emergencystretch=1em
\begin{document}

\maketitle

\begin{abstract}
Стандартный генетический код является почти
универсальным и стабильным на временном масштабе
биологической эволюции, что мотивирует
изучение его арифметических свойств. Мы анализируем числа протонов и нейтронов кодируемых аминокислот (в свободной нейтральной форме) в $33$ группах кодонов, построенных на основе двух естественных систем группировки: разбиения Румера по вырожденности семейств и трёх химических осей третьей позиции кодона (Кето/Амино, Сильные/Слабые, Пурины/Пиримидины). Исключив три
стоп-кодона и инициатор $\mathrm{ATG}$, мы получаем $60$
кодонов, групповые суммы которых определяют арифметическую
структуру по модулю~$37$. Числа нейтронов достигают решётки
с шагом~$37$, тогда как числа протонов~--- неизбежно чётные,
что исключает попадание на нечётные кратные числа~$37$,~---
оказываются на единицу ниже них: как на глобальном уровне,
так и в единственной несущей асимметрию Кето/Амино паре
триптофан/изолейцин, где соответствующие разности равны
$37$ и $36$. Ограничения баланса между ветвями Кето и Амино
вместе с глобальными условиями дают отдельные уравнения
по модулю~$37$ для нейтронов и протонов, причём нейтронное
уравнение содержит $111 = 3 \cdot 37$; поскольку
$\gcd(4, 37) = 1$, эти уравнения однозначно задают
минимальное неотрицательное целочисленное назначение для
четырёх служебных кодонов, после чего регулярности образуют
связное семейство тождеств и делимостей. Среди $27$ таблиц
трансляции NCBI это сочетание обнаруживается только в
стандартном коде и его точном эквиваленте по соответствию
кодонов аминокислотам; детерминированные контроли
показывают, что делимости сосредоточены в группировках по
третьей позиции и что среди проверенных простых модулей
выделяется $37$. В рамках нулевой модели случайного кода наблюдаемые значения ряда арифметических характеристик этой структуры попадают в дальние хвосты соответствующих распределений, полученных методом Монте-Карло. При этом регулярности являются свойствами сумм
по группам кодонов, а не отдельных аминокислот, и зависят
совместно от вырожденности кода и нуклонного состава.
\end{abstract}

\tableofcontents

\section{Введение}
\label{sec:introduction}

Генетический код является одной из наиболее
древних и консервативных молекулярных
характеристик клеточных организмов. Его почти
полная универсальность в известных доменах
жизни и ограниченный масштаб вариаций в
известных альтернативных кодах указывают на
то, что основная структура стандартного кода
уже была установлена на ранних стадиях
биологической эволюции. Поэтому любые организационные закономерности, присутствующие в его арифметической структуре, требуют точного анализа как один из немногих потенциальных источников информации о ранних стадиях жизни.

Арифметическая структура стандартного генетического кода начала изучаться вскоре после его расшифровки. Ю.~Б.~Румер~\cite{rumer1966,rumer1968}
ввёл разбиение $64$ кодонов на два класса по
$32$ кодона, Октет~I
(восемь четырёхкратно вырожденных семейств) и
Октет~II
(остальные кодоны), связанные инволюцией
$R = (\mathrm{T} \leftrightarrow
\mathrm{G}, \mathrm{C} \leftrightarrow
\mathrm{A})$. Щербак и
Макуков~\cite{makukov2013} заметили, что
число $37$ играет повторяющуюся роль в
арифметике кода, появляясь как делитель
нуклонных сумм аминокислот для различных
группировок кодонов. Панов и
Филатов~\cite{panov2024} пересмотрели этот
анализ, введя правило подсчёта по кодонам,
которое образует методологическую основу
настоящего исследования.

Настоящий анализ использует числа протонов и нейтронов двадцати канонических аминокислот в свободной нейтральной форме в рамках соглашения о наиболее распространённых изотопах; после фиксации этих соглашений числа определяются однозначно. Эти величины рассматриваются вместе с двумя естественными системами группировки, заданными независимо от нуклонной арифметики.
Первая представляет собой разбиение Румера.
Вторая задаётся тремя химическими осями
третьей позиции кодона ---
Кето/Амино, Сильные/Слабые и
Пурины/Пиримидины, --- которые исчерпывают
все разбиения четырёх нуклеотидных оснований
на две пары и соответствуют дихотомиям
функциональных групп, числа водородных
связей и топологии колец в химии
нуклеооснований.

Настоящая работа показывает, что в рамках
этой системы арифметические закономерности
стандартного генетического кода организуются
в связное семейство тождеств и
делимостей. Анализ сводит эту связную
структуру к небольшому числу эмпирических
свойств стандартного кода. Затем одни и те же арифметические критерии применяются ко всем $27$
таблицам трансляции NCBI. Среди них проверяемое сочетание структурных
закономерностей обнаруживается только в таблице стандартного
генетического кода и в её точном эквиваленте по соответствию кодонов
аминокислотам. Специфичность
структуры дополнительно исследуется
детерминированными контролями её зависимости от третьей позиции кодона, от модуля и от молекулярного представления аминокислот, а также методом
Монте-Карло относительно явной нулевой модели
случайных кодов.

Анализ разделён на три уровня. Основные закономерности устанавливаются на независимом от параметризации уровне, на котором три стоп-кодона и $\mathrm{ATG}$ в его инициаторной роли исключаются, а суммируются только оставшиеся $60$ кодонов; это центральные результаты работы, и они не предполагают никаких присвоенных значений. Однако чтобы получить суммы для полной таблицы из $64$ кодонов, четырём служебным кодонам необходимо приписать какие-либо нуклонные вклады. Второй уровень --- это наивное прочтение (Key~0), при котором стоп-кодонам приписывается ноль, а $\mathrm{ATG}$ учитывается как кодируемый им метионин. Приписывание нуля также является выбором значения, а не его
отсутствием. Третий уровень~--- это симметризация (Key~1), при которой
четыре служебных кодона несут выведенные нуклонные вклады. Key~1 не
устанавливает независимое от параметризации ядро структуры: оно уже
обнаруживается в суммах по $60$ кодонам, остающимся после исключения
служебных кодонов. Задача Key~1~--- показать, как эта структура
распространяется на полную таблицу с учётом распределения служебных
кодонов по группам, которое задано стандартным кодом. При выбранных условиях
симметрии эти вклады определяются как единственное минимальное
неотрицательное целочисленное решение соответствующей системы
(раздел~\ref{sec:results}), а возникающие закономерности являются
следствиями этого решения.

Приписывание значений служебным кодонам~--- формальный приём,
позволяющий рассматривать арифметику полной таблицы как единый объект;
ни эти значения, ни содержащие их суммы не получают физической
интерпретации.

\section{Методы}
\label{sec:methods}

В данном разделе определяются нуклонные параметры, систематика
кодонных групп и параметризации служебных кодонов, используемые
на протяжении всего анализа, а также описываются расчётные
наборы данных, на которые он опирается.

\subsection{Нуклонный состав аминокислот}
\label{sec:composition}

На протяжении всей работы каждая из двадцати канонических аминокислот
рассматривается как свободная нейтральная молекула аминокислоты (а не
как остаток в составе пептидной цепи), а её молекулярная формула
берётся из базы соединений \textsc{PubChem}~\cite{kim2023}. Из этой
формулы мы вычисляем четыре целочисленные величины, характеризующие
нуклонное содержание молекулы. Далее $Q(a)$ обозначает любую из
четырёх величин $T$, $P$, $N$ или $\Delta$ для одной молекулы
аминокислоты~$a$. Чувствительность результатов к выбранному
молекулярному представлению исследуется в
разделе~\ref{sec:discussion-controls}.

Число протонов $P$ определяется из элементного состава молекулярной
формулы как
\[
P = \sum_E n_E Z_E,
\]
где $n_E$~--- число атомов элемента $E$ в формуле, а $Z_E$~--- его
атомный номер. Используются атомные номера
$Z(\mathrm{H}) = 1$, $Z(\mathrm{C}) = 6$, $Z(\mathrm{N}) = 7$,
$Z(\mathrm{O}) = 8$ и $Z(\mathrm{S}) = 16$.

Число нейтронов $N$ вычисляется путём сопоставления каждому элементу
его наиболее распространённого природного стабильного изотопа согласно
земным изотопным составам, оценённым IUPAC~\cite{berglund2011}.
Соответственно, мы используем $^{1}\mathrm{H}$,
$^{12}\mathrm{C}$, $^{14}\mathrm{N}$, $^{16}\mathrm{O}$ и
$^{32}\mathrm{S}$. Полное число нейтронов равно
\[
N = \sum_E n_E (A_E - Z_E),
\]
где $A_E$~--- массовое число выбранного изотопа элемента $E$.

Полное число нуклонов $T$ и протон-нейтронная разность $\Delta$
определяются как
\[
T = P + N, \qquad \Delta = P - N.
\]
Обе величины по построению целочисленны. В указанных выше изотопных
предположениях только водород вносит вклад в $\Delta$, поскольку у
всех остальных используемых изотопов число протонов равно числу
нейтронов. Следовательно, $\Delta$ численно равно числу атомов водорода в
молекуле и, в частности, $\Delta \geq 0$ для каждой из двадцати
канонических аминокислот. По линейности то же тождество выполняется
для любой суммы вкладов канонических аминокислот и, в частности, для
всех групповых сумм по пулу без служебных кодонов
(раздел~\ref{sec:genetic-code}).

Чтобы избежать коллизии обозначений, мы используем заглавный символ
$\Delta$ исключительно для протон-нейтронного дисбаланса отдельной
молекулы или кодонной группы, а строчный символ $\delta$ резервируем
за попарными разностями молекулярных величин двух аминокислот (например,
$\delta P = P_i - P_j$ между двумя аминокислотами $i$ и $j$).
Если обе аминокислоты указаны явно, их обозначения записываются в
скобках:
\[
\delta Q(a,b)=Q(a)-Q(b).
\]

Подчеркнём, что $P$, $N$, $T$ и $\Delta$ не являются измеренными
молекулярными массами. Это точные целочисленные величины, вычисляемые
из элементного состава, а все арифметические тождества, приводимые в
настоящей работе, являются точными целочисленными соотношениями, а не
артефактами округления приближённых молекулярных масс.

Нуклонные значения всех двадцати канонических аминокислот собраны во
входном наборе данных \texttt{amino\_acids\_nucleons.csv}, где одна
строка соответствует одной аминокислоте, а столбцы содержат $P$, $N$,
$T$ и $\Delta$. Из этой таблицы вычисляются три сопутствующих набора
данных попарных абсолютных разностей:
\texttt{amino\_acids\_proton\_differences.csv},
\texttt{amino\_acids\_neutron\_differences.csv} и
\texttt{amino\_acids\_nucleon\_differences.csv}. Для любой пары
аминокислот $(i, j)$ мы пишем $\delta P = P_i - P_j$ и аналогично для
$N$ и $T$. Каждый файл группирует все $\binom{20}{2} = 190$ неупорядоченных пар
$(i,j)$ по значению соответствующей абсолютной разности
$|\delta P(i,j)|$, $|\delta N(i,j)|$ или $|\delta T(i,j)|$.

Все дальнейшие вычисления нуклонных величин для кодонных групп
сводятся, по линейности, к суммированию значений для отдельных
аминокислот, табулированных в
\texttt{amino\_acids\_nucleons.csv}, с учётом кратности
соответствующих кодонов. Если группа включает служебные кодоны, их
вклады определяются выбранной параметризацией
(см.\ раздел~\ref{sec:parametrization}).

\subsection{Структура генетического кода}
\label{sec:genetic-code}

Стандартный генетический код задаёт отображение $64$ кодонов
четырёхбуквенного нуклеотидного алфавита
$\{\mathrm{A}, \mathrm{T}, \mathrm{G}, \mathrm{C}\}$ на двадцать
канонических аминокислот и сигнал терминации трансляции. В настоящей
работе мы используем ДНК-нотацию (с тимином $\mathrm{T}$, а не с
урацилом $\mathrm{U}$) и придерживаемся соответствий кодонов
стандартного кода (таблица трансляции
NCBI~1)~\cite{ncbi2024geneticcodes}. Три кодона $\mathrm{TAA}$,
$\mathrm{TAG}$ и $\mathrm{TGA}$ служат стоп-сигналами, а кодон
$\mathrm{ATG}$ кодирует метионин и служит каноническим сигналом
инициации трансляции. Для целей настоящего арифметического анализа
именно эти четыре кодона --- три стоп-кодона и $\mathrm{ATG}$ в его
инициаторной роли --- называются далее \emph{служебными кодонами}.
Среди возможных инициаторных кодонов $\mathrm{ATG}$ является
преобладающим. Его РНК-эквивалент $\mathrm{AUG}$ --- единственный
из них, полностью комплементарный антикодону $\mathrm{CAU}$
инициаторной тРНК метионина. В настоящей работе именно
$\mathrm{ATG}$ представляет функцию инициации; это аналитическое
соглашение не исключает инициации с других кодонов.

Обозначим множество всех $64$ кодонов через
\[
\mathcal{C}_{64}
=
\{\mathrm{A},\mathrm{T},\mathrm{G},\mathrm{C}\}^{3},
\]
а множество четырёх служебных кодонов через
\[
\mathcal{S}
=
\{\mathrm{ATG},\mathrm{TAA},\mathrm{TAG},\mathrm{TGA}\}.
\]
Тогда \emph{60-кодонный пул без служебных кодонов} определяется как
\[
\mathcal{C}_{\mathrm{SE}}
=
\mathcal{C}_{64}\setminus\mathcal{S},
\]
и $|\mathcal{C}_{\mathrm{SE}}|=60$. Все кодоны, принадлежащие
$\mathcal{C}_{\mathrm{SE}}$, учитываются согласно их стандартным
аминокислотным назначениям. Поскольку $\mathrm{ATG}$ является единственным кодоном метионина,
$\mathcal{C}_{\mathrm{SE}}$ содержит кодоны девятнадцати канонических
аминокислот; метионин в этом пуле отсутствует. Полное отображение кодон--продукт
приведено в наборе данных
\texttt{genetic\_code\_codons.csv}.

Кодонным семейством мы называем множество четырёх кодонов с
одинаковыми основаниями в первых двух позициях и произвольным
основанием в третьей позиции, обозначаемым символом $\mathrm{N}$.
Структурным свойством стандартного кода, впервые выявленным
Румером~\cite{rumer1966,rumer1968}, является разбиение $64$ кодонов
на две половины по $32$ кодона в каждой, различающиеся характером
вырожденности третьей позиции кодона. Первая половина, которую мы
обозначаем \emph{Октет~I}, состоит из восьми кодонных семейств,
в которых все четыре нуклеотида в третьей позиции кодируют одну
и ту же аминокислоту. Эти полностью четырёхкратно вырожденные
семейства --- $\mathrm{GCN}$ (Ala), $\mathrm{CGN}$ (Arg),
$\mathrm{GGN}$ (Gly), $\mathrm{CCN}$ (Pro), $\mathrm{ACN}$ (Thr),
$\mathrm{GTN}$ (Val), $\mathrm{CTN}$ (Leu) и $\mathrm{TCN}$ (Ser).
Вторая половина, обозначаемая \emph{Октет~II}, состоит из оставшихся
восьми кодонных семейств, в которых нуклеотид в третьей позиции
влияет на кодируемый продукт. Каждый служебный кодон стандартного
кода принадлежит Октету~II и несёт пурин ($\mathrm{A}$ или
$\mathrm{G}$) в третьей позиции.

Разбиение Румера не является аналитической конструкцией, введённой
в настоящей работе, а представляет собой структурную особенность
самого кода. Оно отражает как качественное различие между полностью
и частично вырожденными кодонными семействами, так и количественный
факт равной мощности этих двух классов: каждый из них содержит
восемь семейств и тридцать два кодона. Поскольку Октет~I не содержит
служебных кодонов, все нуклонные величины, вычисленные по Октету~I,
не зависят от параметризации старт- и стоп-сигналов и определяются
только нуклонным содержанием кодируемых аминокислот. То же свойство
выполняется по отдельной причине для любой кодонной группы, у которой
нуклеотид в третьей позиции ограничен пиримидинами (цитозином или
тимином), поскольку все четыре служебных кодона стандартного кода
несут в третьей позиции пурин. Полная четырёхкратная вырожденность
и отсутствие пуринов в третьей позиции представляют собой два
различных структурных условия, но приводят к одному следствию:
нуклонные величины соответствующих групп фиксированы их
аминокислотным составом и дают не зависящие от параметризации
опорные значения для арифметических тождеств
раздела~\ref{sec:results}.

Ещё одно структурное свойство разбиения Румера, введённое
Румером~\cite{rumer1968}, состоит в том, что два октета связаны
простой перестановкой нуклеотидного алфавита. При преобразовании
\[
R =
(\mathrm{T} \leftrightarrow \mathrm{G},\;
 \mathrm{C} \leftrightarrow \mathrm{A})
\]
каждое кодонное семейство Октета~I отображается в кодонное семейство
Октета~II и наоборот. Это преобразование, известное как
преобразование Румера, перестраивает разбиение, но не сохраняет
аминокислоты, кодируемые отдельными кодонами.

В совокупности четырёхкратная вырожденность и преобразование Румера
выделяют разбиение на Октет~I / Октет~II как естественное структурное
деление стандартного кода, независимое от нуклонной арифметики.
В настоящей работе это деление используется как опорное разбиение,
по которому вычисляются и сопоставляются нуклонные величины.

Таблица кодонов (Рис.~\ref{fig:codon-table}) отображается с осями
в порядке оснований $C,G,T,A$, который также является порядком по
умолчанию в сопровождающем визуализаторе. Этот порядок имеет два
обоснования. С химической точки зрения он был введён
Румером~\cite{rumer1968} на основе типа основания и числа
водородных связей в канонической паре оснований.
Последовательность $C,G,T,A$ чередует пиримидины и пурины,
а внутри каждого из этих классов основание из пары Сильных
$\{\mathrm{C},\mathrm{G}\}$ предшествует основанию из пары Слабых
$\{\mathrm{T},\mathrm{A}\}$: $\mathrm{C}$ предшествует $\mathrm{T}$
среди пиримидинов, а $\mathrm{G}$ предшествует $\mathrm{A}$ среди
пуринов. С геометрической точки зрения, как было показано Пановым
и Филатовым~\cite{panov2024}, это один из ровно четырёх порядков
(из $24$ возможных перестановок четырёх оснований), при которых
Октет~I и Октет~II
образуют связные области на матрице кодонов.

Чтобы оценить, в какой мере изучаемые арифметические свойства
специфичны для стандартного генетического кода, мы дополнительно
извлекли полный набор из $27$ таблиц трансляции, поддерживаемых
NCBI~\cite{ncbi2024geneticcodes}. Таблицы имеют номера 1--33,
причём часть номеров не используется, и охватывают ядерные,
митохондриальные и пластидные коды в широком таксономическом
диапазоне. Соответствия кодонов во всех этих таблицах собраны
в наборе данных
\texttt{ncbi\_genetic\_code\_registry.csv} и используются в
разделе~\ref{sec:results-key1-ncbi} для сопоставления стандартного
кода с остальными таблицами NCBI.

\begin{figure}[!htbp]
  \centering
  \includegraphics[width=0.85\textwidth]{codon_table_key0.png}
\caption{Стандартный генетический код с нуклонными параметрами
кодируемых аминокислот, отображённый при наивном прочтении
(Key~0), при котором стоп-кодоны имеют нулевой вклад, а
$\mathrm{ATG}$ учитывается как кодируемый им метионин. В каждой
ячейке кодон расположен слева; обозначение кодируемого продукта
(трёхбуквенный символ аминокислоты либо $\mathrm{STOP}$ для трёх
стоп-кодонов) --- в верхнем центре; полное число нуклонов $T$ ---
под символом продукта; число протонов $P$ --- в правом верхнем
углу; число нейтронов $N$ --- в правом нижнем углу. Кодоны
Октета~I (восемь полностью четырёхкратно
вырожденных семейств) выделены фиолетовым, а кодоны
Октета~II --- жёлтым. Сгенерировано
сопровождающим \texttt{visualization.html}.}
  \label{fig:codon-table}
\end{figure}

\subsection{Построение кодонных групп}
\label{sec:partitioning}

В дополнение к разбиению Румера, введённому в
разделе~\ref{sec:genetic-code}, мы используем дополнительную
перекрывающуюся систему группировки, основанную на химической
идентичности нуклеотида в третьей позиции кодона. Каждый из четырёх
нуклеотидов относится к одному из двух классов в каждой из трёх
различных бинарных химических классификаций:
\begin{itemize}
  \item тип основания: \emph{пурин}
        $\{\mathrm{A},\mathrm{G}\}$ (R) или \emph{пиримидин}
        $\{\mathrm{C},\mathrm{T}\}$ (Y);
  \item число водородных связей в канонической паре оснований:
        \emph{сильная} пара $\{\mathrm{C},\mathrm{G}\}$ (S; три
        связи) или \emph{слабая} пара
        $\{\mathrm{A},\mathrm{T}\}$ (W; две связи);
  \item функциональный класс: \emph{амино}
        $\{\mathrm{A},\mathrm{C}\}$ (M) или \emph{кето}
        $\{\mathrm{G},\mathrm{T}\}$ (K).
\end{itemize}
Эти оси не были выбраны по результатам последующих нуклонных
вычислений: они представляют собой три общепринятые парные
дихотомии нуклеотидов и явно закреплены в номенклатуре
неоднозначности IUPAC--IUB парами символов
R/Y, S/W и M/K. Вместе они исчерпывают все три способа
разбить множество из четырёх нуклеотидов на две
неупорядоченные пары, что даёт шесть двухэлементных
подмножеств
$X \subset \{\mathrm{A}, \mathrm{T}, \mathrm{G}, \mathrm{C}\}$.

Для каждого из этих шести двухэлементных подмножеств, а также для
каждого одноэлементного подмножества
$X\in\{\{\mathrm{A}\},\{\mathrm{C}\},\{\mathrm{G}\},
\{\mathrm{T}\}\}$, определим кодонную группу
\[
G_X =
\{(c_1,c_2,c_3)\in\mathcal{C}_{64}
  : c_3\in X\}.
\]
Иными словами, $G_X$ состоит из всех кодонов, у которых нуклеотид
в третьей позиции принадлежит $X$. Шесть двухнуклеотидных групп
перекрываются: каждый кодон входит ровно в три из них, по одной на
каждую химическую классификацию. Одновременно каждый кодон входит
ровно в одну из четырёх однонуклеотидных групп.

Эта система группировки применяется на трёх уровнях: ко всем
$64$ кодонам, к $32$ кодонам Октета~I и к $32$ кодонам
Октета~II. Если $|X|=2$, то $G_X$ содержит $32$ кодона, а каждое
из пересечений $G_X\cap\mathrm{Octet~I}$ и
$G_X\cap\mathrm{Octet~II}$ содержит по $16$ кодонов. Если
$|X|=1$, соответствующие размеры равны $16$ и $8$ кодонам.
Размер сам по себе не идентифицирует конкретную группу, поэтому
в формулах мы указываем множество $X$ и, при необходимости,
октет явно.

Используемое в настоящей работе семейство из $33$ кодонных групп
включает полный код как одну группу из $64$ кодонов; два октета
как группы по $32$ кодона; десять групп $G_X$ на уровне полного
кода, соответствующих шести двухнуклеотидным и четырём
однонуклеотидным множествам $X$; и двадцать их пересечений с
Октетом~I и Октетом~II. Таким образом, общее число анализируемых
групп равно
\[
1+2+10+20=33.
\]
Например, $G_{\{\mathrm{G},\mathrm{T}\}}\cap
\mathrm{Octet~I}$ представляет группу
$\mathrm{Keto}\cap\mathrm{Octet~I}$ из $16$ кодонов, а
$G_{\{\mathrm{A}\}}\cap\mathrm{Octet~II}$ --- группу из $8$
кодонов Октета~II с аденином в третьей позиции. Все $33$ группы
вместе с их размерами, кодонными составами и принадлежностью к
октетам Румера перечислены в наборе данных
\texttt{codon\_groups.csv}.

В реестре групп полному множеству $\mathcal{C}_{64}$ соответствует
строка с id~1. В статье для этой группы используется обозначение
$\mathrm{All}$, тогда как $\mathcal{C}_{64}$ обозначает само множество
кодонов.

\subsection{Пиримидиновая синонимия и тождества групп-близнецов}
\label{sec:twin-identities}

Для каждого кодонного семейства стандартного генетического кода кодоны
$\mathrm{NNC}$ и $\mathrm{NNT}$ с одинаковыми первыми двумя позициями
кодируют одну и ту же аминокислоту. Поэтому замена
$\mathrm{C}\leftrightarrow\mathrm{T}$ в третьей позиции не меняет
мультимножество кодируемых аминокислот. Она переводит Кето
$\{\mathrm{G},\mathrm{T}\}$ в Сильные
$\{\mathrm{C},\mathrm{G}\}$, а Амино
$\{\mathrm{A},\mathrm{C}\}$~--- в Слабые
$\{\mathrm{A},\mathrm{T}\}$.

Члены каждой пары содержат одни и те же служебные кодоны: Кето и
Сильные содержат $\mathrm{ATG}$ и $\mathrm{TAG}$, а Амино и
Слабые~--- $\mathrm{TAA}$ и $\mathrm{TGA}$. Поэтому при любом
назначении $K$, сопоставляющем каждому служебному кодону пару значений
$(P,N)$ (раздел~\ref{sec:parametrization}), и для
$Q\in\{T,P,N,\Delta\}$ выполняются тождества
\[
Q_K(\mathrm{Keto})=Q_K(\mathrm{Strong}),\qquad
Q_K(\mathrm{Amino})=Q_K(\mathrm{Weak}).
\]
Те же тождества сохраняются после пересечения обеих групп каждой пары
с одним и тем же октетом. Равенства не зависят от параметризации, хотя
сами значения полных групп могут зависеть от~$K$. Пары Кето/Сильные и
Амино/Слабые, а также соответствующие пары после пересечения с одним и
тем же октетом мы называем \emph{группами-близнецами}. В наборе данных
\texttt{codon\_groups.csv} это отношение фиксируется в столбце
\texttt{Twin\_Group\_ID}.

Пурины и Пиримидины не являются группами-близнецами: замена
$\mathrm{C}\leftrightarrow\mathrm{T}$ сохраняет каждую из этих групп,
а не переводит одну в другую. Та же пиримидиновая синонимия
непосредственно даёт равенство всех четырёх нуклонных величин для
однонуклеотидных групп $\{\mathrm{C}\}$ и $\{\mathrm{T}\}$, а также
для их пересечений с каждым октетом. Эти группы и их пересечения с
октетами не содержат служебных кодонов, поэтому как сами значения, так
и равенства между ними не зависят от параметризации. Поскольку это
пара однонуклеотидных групп, а не пара двухнуклеотидных ветвей
химических осей, термин \emph{группы-близнецы} и столбец
\texttt{Twin\_Group\_ID} для неё не используются.

\subsection{Параметризация служебных кодонов}
\label{sec:parametrization}

Нуклонные величины кодонной группы получают суммированием значений
аминокислот, кодируемых входящими в неё кодонами стандартного
генетического кода. Для кодонов, учитываемых только по их
аминокислотному назначению, такое суммирование однозначно. Отдельного
соглашения требуют три стоп-кодона, которым не соответствует
аминокислотный продукт, и $\mathrm{ATG}$ в его инициаторной роли;
сам $\mathrm{ATG}$ также кодирует метионин
(раздел~\ref{sec:genetic-code}).

Множество всех $64$ кодонов $\mathcal{C}_{64}$, множество четырёх
служебных кодонов $\mathcal{S}$ и остающийся 60-кодонный пул
$\mathcal{C}_{\mathrm{SE}}$ определены в
разделе~\ref{sec:genetic-code}. Далее $\mathcal{C}_{\mathrm{SE}}$
называется \emph{пулом без служебных кодонов}; в коде и именах файлов
ему отвечает идентификатор \texttt{pool60}.
Каждый кодон из $\mathcal{C}_{\mathrm{SE}}$ однозначно сопоставлен
канонической аминокислоте. Для
$c\in\mathcal{C}_{\mathrm{SE}}$ и
$Q\in\{T,P,N,\Delta\}$ через $Q(c)$ обозначается соответствующая
величина свободной нейтральной формы аминокислоты, кодируемой кодоном
$c$, согласно разделу~\ref{sec:composition}.

Для любого множества кодонов
$A\subseteq\mathcal{C}_{64}$, не содержащего служебных кодонов,
положим
\[
Q(A):=\sum_{c\in A}Q(c).
\]
Для группы кодонов $G$ её сумма по пулу без служебных кодонов
определяется как
\[
Q_{\mathrm{SE}}(G)
:=
\sum_{c\in G\cap\mathcal{C}_{\mathrm{SE}}}Q(c).
\]
Для служебного кодона $c$ обозначение $Q(c)$ без индекса ключа не
используется, поскольку формальный вклад такого кодона зависит от
выбранного ключа. Индекс
$\mathrm{SE}$ обозначает ограничение области суммирования, а не
параметрический ключ. Ни сумма $Q_{\mathrm{SE}}(G)$, ни число входящих
в неё кодонов $|G\cap\mathcal{C}_{\mathrm{SE}}|$ не зависят от
назначений служебным кодонам.

Если сама группа $G$ не содержит служебных кодонов, то ограничение
ничего не исключает и
\[
Q(G)=Q_{\mathrm{SE}}(G).
\]

Для анализа полной таблицы из $64$ позиций введём
\emph{параметрический ключ} $K$: фиксированное назначение пары чисел
$\bigl(P_K(c),N_K(c)\bigr)$ каждому служебному кодону $c$.
При таком ключе групповые суммы определяются как
\[
\bigl(P_K(G),N_K(G)\bigr)
:=
\bigl(P_{\mathrm{SE}}(G),N_{\mathrm{SE}}(G)\bigr)
+
\sum_{c\in G\cap\mathcal{S}}
\bigl(P_K(c),N_K(c)\bigr),
\]
а две остальные величины определяются, а не назначаются независимо:
\[
T_K(G)=P_K(G)+N_K(G),
\qquad
\Delta_K(G)=P_K(G)-N_K(G).
\]
Те же два соотношения задают $T_K(c)$ и $\Delta_K(c)$ для каждого
служебного кодона~$c$.
Для $Q\in\{T,P,N,\Delta\}$ через $Q_0(G)$ и $Q_1(G)$ далее
обозначаются полные групповые суммы, включая вклады служебных кодонов,
соответственно при назначениях Key~0 и Key~1. Индексы $0$ и $1$
являются метками этих двух назначений, а не числовыми значениями
скалярного параметра. Таким образом, $Q_{\mathrm{SE}}(G)$ не зависит
от ключа, а $Q_0(G)$ и $Q_1(G)$ могут различаться только вкладами
служебных кодонов, содержащихся в группе $G$. Если группа $G$ не
содержит служебных кодонов, то
$Q_K(G)=Q(G)=Q_{\mathrm{SE}}(G)$ при любом $K$. Индекс
$\mathrm{SE}$ применяется только к группам кодонов; индексы $K$, $0$
и $1$~--- к полным групповым суммам и формальным вкладам отдельных
служебных кодонов. Два рассматриваемых ключа зафиксированы в
\texttt{key\_parameters.csv}.

\paragraph{Наивное прочтение (Key~0).}
Каждому из трёх стоп-кодонов приписывается нулевой вклад
$(P,N)=(0,0)$, поскольку стоп-сигнал не имеет аминокислотного продукта.
Кодон $\mathrm{ATG}$ трактуется как обычный кодон метионина и получает
значения свободной молекулы метионина,
$(P,N)=(80,69)$. Это простое соглашение соответствует обычному чтению
аминокислот-кодирующей части таблицы, в котором стоп-сигналам
приписан нулевой вклад, но само остаётся параметризацией. В частности,
Key~0 не тождествен пулу без служебных кодонов: Key~0 сохраняет все
$64$ позиции и учитывает $\mathrm{ATG}$ как метионин, тогда как в пуле
суммирование идёт только по оставшимся $60$ кодонам.

\paragraph{Симметризация (Key~1).}
Каждому из трёх стоп-кодонов приписывается
$(P,N)=(37,37)$, а $\mathrm{ATG}$ как формальному маркеру инициации~---
$(P,N)=(1,0)$. Эти значения не интерпретируются как нуклонный состав
физических продуктов трансляции. Они выводятся в
разделах~\ref{sec:results-key1-derivation}--%
\ref{sec:results-key1-uniqueness} в предположении, что три стоп-кодона
имеют общую пару параметров, и при наложении указанных там условий
баланса и делимости на~$37$. При этих предпосылках Key~1 является
единственным минимальным неотрицательным целочисленным решением
полученной системы. Минимальность здесь означает наименьший суммарный вклад
служебных кодонов~--- сумму вкладов трёх стоп-кодонов и
$\mathrm{ATG}$,~--- определяемый отдельно для протонов и
для нейтронов. Для стандартного кода баланс фиксирует
разность параметров стоп-кодона и инициатора, поэтому
суммарный вклад возрастает с параметром инициатора и
минимален при его наименьшем допустимом значении. В настоящем разделе Key~1 вводится
как результат этой формальной конструкции, а не как биохимическое
свойство стоп-сигналов или инициатора.

Для кодонных групп, не содержащих служебных кодонов, значения при
Key~0 и Key~1 совпадают с непосредственными аминокислотными суммами.
Если группа содержит служебные кодоны, назначения Key~1 имеют только
формальный арифметический смысл: соответствующие суммы не являются
нуклонным составом физического ансамбля молекул. Обозначения $P$, $N$, $T$ и $\Delta$ сохраняются
для единообразия с исходными аминокислотными вкладами,
но биохимическая интерпретация формальных вкладов Key~1
не утверждается.

Набор \texttt{key\_parameters.csv} дополнительно содержит восемь
справочных назначений \texttt{Key~2}--\texttt{Key~9}, рассмотренных на
исследовательском этапе. Они не используются для получения результатов
статьи и не обладают статусом выводимого Key~1. Далее три аналитических
уровня всегда различаются явно: пул без служебных кодонов, наивное
прочтение (Key~0) и симметризация (Key~1).

\subsection{Аналитические критерии и классификация равенств}
\label{sec:criteria}

Для каждой из $33$ кодонных групп раздела~\ref{sec:partitioning} мы
рассматриваем величины $P$, $N$, $T$ и $\Delta$ на трёх аналитических
уровнях: в пуле без служебных кодонов, при Key~0 и при Key~1. Ниже
описаны три процедуры регистрации результатов:
проверка делимости на~$37$ и десятичной круглости, а также поиск точных
равенств. Делимость на~$37$ является основным заранее выбранным
арифметическим критерием. Точные степени, сокращённые отношения и
совпадения с разностями аминокислот рассматриваются отдельно как
дополнительные описательные наблюдения. Статистическая значимость
оценивается только там, где задана соответствующая нулевая модель.

\paragraph{Критерий 1: делимость на~37.}
Для каждой групповой величины $P$, $N$, $T$ и $\Delta$ вычисляются
частное и остаток при делении на~$37$. Нулевой остаток означает точную
делимость. Этот критерий продолжает направление исследований,
начатое \citet{makukov2013} и развитое \citet{panov2024}. Если значение
делится на~$37$, разложение также показывает дальнейшие множители
частного; им не приписывается самостоятельный доказательный вес без
отдельного анализа. Ненулевые остатки сообщаются описательно лишь тогда, когда они
входят в точные соотношения между заранее определёнными группами или
образуют повторяющиеся паттерны единичных отклонений от ближайшего
кратного~$37$.

Полные выражения вида $q\cdot37+r$ для всех групп и четырёх величин
приведены в
\texttt{pool60\_divisibility-37.csv},
\texttt{key0\_divisibility-37.csv} и
\texttt{key1\_divisibility-37.csv}.

\paragraph{Критерий 2: круглость в десятичной системе.}
Положительную целую величину мы называем круглой в десятичной системе,
если она делится на $10^k$ хотя бы для одного целого $k\geq1$; степень
круглости задаётся наибольшим таким $k$. Этот критерий используется
описательно, поскольку в данных встречаются точные кратные $100$ и
$1000$, в частности среди не зависящих от параметризации сумм
Октета~I. Хотя понятие круглости связано с десятичной записью, проверяемая
делимость является свойством самого целого числа: условие
$10^k\mid x$ эквивалентно одновременной делимости на $2^k$ и $5^k$.
Систематический контроль по простым модулям приведён в
разделе~\ref{sec:discussion-controls}.

\paragraph{Критерий 3: точные равенства.}
На каждом аналитическом уровне мы сравниваем друг с другом значения
$P$, $N$, $T$ и $\Delta$ для всех $33$ групп и регистрируем все точные
совпадения. Тривиальные тождества вида $P(G)=P(G)$ в результаты не
включаются. Каждое найденное равенство относится к одному из трёх
классов:
\begin{itemize}
  \item \emph{равенства между группами разных размеров}
        (в столбце \texttt{Match\_Type}~--- \texttt{DIFF\_PARAM}):
        одинаковое численное значение получается при суммировании по
        разному числу кодонов; сравниваться при этом могут как
        одинаковые, так и разные нуклонные величины;
  \item \emph{межвеличинные равенства групп одного размера}
        (\texttt{CROSS\_PARAM}): группы содержат одинаковое число
        кодонов, но сравниваются разные величины, например
        $P(G_1)=N(G_2)$;
  \item \emph{равенства одной величины для групп одного размера}
        (\texttt{SAME\_PARAM}): группы содержат одинаковое число
        кодонов и сравнивается одна и та же величина, например
        $P(G_1)=P(G_2)$.
\end{itemize}

Для классификации используется число кодонов, фактически включённых в
сумму. При Key~0 и Key~1 это полный размер группы в таблице из $64$
кодонов. В пуле без служебных кодонов служебные кодоны исключаются,
поэтому содержащие их группы имеют меньший размер.

Все найденные равенства приведены в
\texttt{pool60\_equalities.csv},
\texttt{key0\_equalities.csv} и
\texttt{key1\_equalities.csv}. Тождества групп-близнецов из
раздела~\ref{sec:twin-identities} также регистрируются, но рассматриваются
как следствия одинакового аминокислотного состава, а не как новые
совпадения.

\paragraph{Дополнительные исследовательские наблюдения.}
Помимо основных процедур мы описательно отмечаем:
\begin{itemize}
  \item точные степени малых целых чисел, например
        $4096=2^{12}$;
  \item сокращённые отношения внутри одной группы, включая
        $P/T$, $N/T$, $\Delta/T$, $P/N$, $\Delta/P$ и $\Delta/N$,
        табулированные в
        \texttt{pool60\_ratios.csv},
        \texttt{key0\_ratios.csv} и
        \texttt{key1\_ratios.csv};
  \item совпадения групповых величин с абсолютными разностями
        $|\delta P|$, $|\delta N|$ или $|\delta T|$ отдельных пар
        аминокислот, например $|\delta N(\mathrm{Trp},\mathrm{Ile})|
        =37$.
\end{itemize}
Эти наблюдения описывают арифметику и могут служить основой для
проверяемых гипотез.

Наконец, четыре величины не являются независимыми:
\[
T=P+N,
\qquad
\Delta=P-N.
\]
Поэтому несколько равенств, делимостей или отношений, получаемых друг
из друга этими тождествами, не считаются отдельными свидетельствами.
То же относится к копиям, распространяющимся через вложенность групп
или тождества групп-близнецов. В разделе~\ref{sec:results} мы
различаем: (i) факты, непосредственно вычисленные для пула без
служебных кодонов и не требующие назначений служебным кодонам;
(ii) тождества,
для которых независимость от ключа доказана алгебраически; и
(iii) результаты, специфичные для Key~0 или Key~1.

\subsection{Вычисляемые наборы данных и воспроизводимость}
\label{sec:datasets}

\begingroup\sloppy

Основные детерминированные результаты настоящей работы производятся
скриптом \texttt{reproduce.py}. Он читает пять входных наборов данных
и создаёт $18$ вычисляемых выходных наборов. Отдельный
детерминированный скрипт \texttt{controls.py} создаёт четыре
контрольных набора. Основной статистический расчёт выполняется
скриптом \texttt{mc\_significance.py}, дополнительные проверки~---
скриптом \texttt{mc\_sensitivity.py}. Они создают таблицы
\texttt{mc\_significance.csv} и \texttt{mc\_sensitivity.csv}
соответственно. Результаты независимо проверяются скриптами
\texttt{verify\_mc.py} и \texttt{verify\_mc\_sensitivity.py}.
По умолчанию все пути к данным разрешаются относительно каталога
соответствующего скрипта.

\paragraph{Входные наборы данных.}
Пять входных файлов задают исходный численный и структурный слой
анализа:
\begin{itemize}
  \item \texttt{amino\_acids\_nucleons.csv} содержит значения
        $P$, $N$, $T$ и $\Delta$ двадцати канонических аминокислот
        (раздел~\ref{sec:composition});
  \item \texttt{genetic\_code\_codons.csv} задаёт стандартное
        соответствие кодонов и продуктов трансляции
        (раздел~\ref{sec:genetic-code});
  \item \texttt{ncbi\_genetic\_code\_registry.csv} фиксирует снимок
        $27$ таблиц трансляции NCBI, включая стандартную таблицу
        (раздел~\ref{sec:genetic-code});
  \item \texttt{codon\_groups.csv} определяет систему из $33$
        перекрывающихся кодонных групп
        (раздел~\ref{sec:partitioning});
  \item \texttt{key\_parameters.csv} содержит назначения служебным
        кодонам для параметрических ключей
        (раздел~\ref{sec:parametrization}).
\end{itemize}

\paragraph{Вычисляемые наборы данных по ключам.}
Для каждого из двух основных ключей, Key~0 и Key~1,
\texttt{reproduce.py} создаёт четыре файла. Префиксы
\texttt{key0\_} и \texttt{key1\_} обозначают использованный ключ:
\begin{itemize}
  \item \texttt{*\_nucleon-data.csv} содержит четыре сообщаемые
        групповые величины $P$, $N$, $T$ и $\Delta$ для всех $33$
        групп; $T$ и $\Delta$ вычисляются из $P$ и $N$;
  \item \texttt{*\_divisibility-37.csv} представляет каждую
        величину в точной форме $q\cdot37+r$;
  \item \texttt{*\_equalities.csv} исчерпывающе перечисляет точные
        численные равенства и классифицирует их в столбце
        \texttt{Match\_Type} как \texttt{DIFF\_PARAM},
        \texttt{CROSS\_PARAM} или \texttt{SAME\_PARAM} согласно
        разделу~\ref{sec:criteria};
  \item \texttt{*\_ratios.csv} содержит шесть отношений
        $P/T$, $N/T$, $\Delta/T$, $P/N$, $\Delta/P$ и $\Delta/N$
        в точной сокращённой дробной форме и в виде округлённого
        десятичного приближения.
\end{itemize}
Таким образом, два ключа дают восемь выходных файлов.

\paragraph{Наборы данных пула без служебных кодонов.}
Для пула без служебных кодонов скрипт создаёт четыре аналогичных
файла:
\texttt{pool60\_nucleon-data.csv},
\texttt{pool60\_divisibility-37.csv},
\texttt{pool60\_equalities.csv} и
\texttt{pool60\_ratios.csv}. Перед суммированием из каждой группы $G$
исключаются служебные кодоны, так что число учитываемых кодонов равно
$|G\cap\mathcal{C}_{\mathrm{SE}}|$. Никакой параметрический ключ при
этом не применяется.

\paragraph{Аминокислотные разности.}
Три файла, создаваемые один раз,
\texttt{amino\_acids\_proton\_differences.csv},
\texttt{amino\_acids\_neutron\_differences.csv} и
\texttt{amino\_acids\_nucleon\_differences.csv}, перечисляют для
всех $\binom{20}{2}=190$ неупорядоченных пар канонических
аминокислот абсолютные разности $|\delta P|$, $|\delta N|$ и
$|\delta T|$ соответственно.

\paragraph{Таблично-зависимые пулы NCBI.}
Три оставшихся выхода основного скрипта сравнивают стандартный код
с остальными таблицами в
\texttt{ncbi\_genetic\_code\_registry.csv}. Поскольку
наборы стоп-кодонов различаются между таблицами, мы не предполагаем
заранее, что каждый анализируемый пул содержит $60$ кодонов.
Используются две модели. Модель~1 исключает
$\mathrm{ATG}$ и все кодоны, перечисленные в поле
\texttt{Stop\_Codons}. Модель~2 исключает $\mathrm{ATG}$ и только
постоянные стоп-кодоны, обозначенные символом \texttt{*} в поле
\texttt{Amino\_Acids}; контекстно-зависимые стопы включаются в анализ
с указанными для них аминокислотными назначениями. Для таблиц без
контекстно-зависимых стопов модели совпадают. Для стандартной
таблицы обе модели дают один и тот же 60-кодонный пул без служебных
кодонов. Фактическое число кодонов в
анализируемом пуле записывается для обеих моделей (столбцы
\texttt{Pure\_Codons\_Count} и \texttt{Pure\_Codons\_Count\_Alt}).

\texttt{deficit\_models\_analysis.csv} содержит по одной строке для
каждой из $27$ таблиц. Полные протонная и нейтронная суммы вычисляются
по зарегистрированным аминокислотным
назначениям, причём позиции \texttt{*} вносят нулевой вклад. Для
обеих моделей дополнительно вычисляются суммы по
соответствующему анализируемому пулу, остаток протонной суммы по
модулю~$37$ и три нейтронных дефицита относительно следующих
опорных величин:
\begin{enumerate}
  \item константы $3700$, полученной для Октета~I стандартного кода;
  \item структурного $T(\mathrm{Octet~I})$, где сохраняются восемь
        семейств, занимающих позиции Октета~I стандартной таблицы;
  \item функционального $T(\mathrm{Octet~I})$, где в самой таблице
        $t$ выбираются семейства из четырёх кодонов с одним
        аминокислотным назначением и без постоянного стоп-кодона.
\end{enumerate}

\texttt{keto\_amino\_balance\_models.csv} использует те же две
модели и для каждой таблицы суммирует $P$ и $N$ отдельно по
пересечениям анализируемого пула с группами Кето и Амино третьей
позиции. В файл записываются обе суммы и их направленная разность
Кето минус Амино.

\texttt{service\_codon\_balance\_comparison.csv} содержит
$54$ строки: по одной для каждой таблицы и каждой модели.
Для протонов и нейтронов отдельно определяются
неотрицательные целочисленные назначения с минимальным
суммарным вкладом служебных кодонов. Всем стоп-кодонам
принятой модели приписывается одинаковое значение,
а $\mathrm{ATG}$~--- отдельное. Сравниваются две процедуры:
баланс Кето/Амино и тот же баланс с дополнительным
требованием делимости полной суммы на~$37$.
Файл содержит назначения, добавленные вклады,
итоговые суммы и разность минимальных вкладов между
процедурами. Отдельно различаются единственное назначение,
определённость только суммарного вклада и отсутствие
решения. Применение этих процедур к стандартному коду
рассмотрено в разделе~\ref{sec:results-key1-derivation}.

\paragraph{Воспроизводимость основного конвейера.}
\texttt{reproduce.py} использует только стандартную библиотеку
Python и не требует внешних пакетов или виртуального окружения.
Вычисления детерминированы: выходы записываются в UTF-8 с окончаниями
строк LF и фиксированным порядком строк. Канонические файлы
репозитория воспроизводятся из пяти входов без ручного редактирования.
Все групповые суммы и остатки вычисляются в целых числах, а отношения
сначала сокращаются точно посредством наибольшего общего делителя.
Десятичные столбцы в \texttt{*\_ratios.csv} являются только
округлёнными представлениями точных дробей и не используются для
получения результатов. Каждое детерминированное численное утверждение
можно проследить до конкретной ячейки или документированной точной
комбинации ячеек выходных файлов.

Рандомизация отсутствует как в \texttt{reproduce.py}, так и в
\texttt{controls.py}. Она используется в основном расчёте
Монте-Карло \texttt{mc\_significance.py}, дополнительных проверках
\texttt{mc\_sensitivity.py} и их независимых реализациях
\texttt{verify\_mc.py} и \texttt{verify\_mc\_sensitivity.py}.
Во всех случаях заданы начальные значения генератора и число
испытаний. Оценки хвостовых вероятностей и их интервалы
представлены числами с плавающей точкой; принадлежность хвосту
определяется без округления показателей. Статистические процедуры
описаны в разделе~\ref{sec:discussion-montecarlo}.

\paragraph{Вспомогательный визуализатор.}
Файл \texttt{visualization.html} предоставляет интерактивное
представление таблиц и кодонных групп в режимах пула без служебных
кодонов (POOL60), Key~0 и Key~1, а также дополнительные
исследовательские представления. Он не входит в основной
воспроизводимый конвейер и не является источником опубликованных
чисел: детерминированные численные значения основного текста
проверяются по соответствующим CSV. Онлайн-версия доступна по адресу
\url{https://ruslankhafizov.github.io/genetic-code-arithmetic/}.

\paragraph{Вспомогательные контрольные наборы данных.}
\texttt{controls.py} читает те же пять входных файлов, использует
только стандартную библиотеку Python и создаёт четыре
детерминированных набора:
\begin{itemize}
  \item \texttt{position\_axis\_analysis.csv} применяет три
        бинарные химические классификации Кето/Амино,
        Сильные/Слабые и Пурины/Пиримидины по отдельности для первой,
        второй и третьей позиций кодона на трёх аналитических
        уровнях: в пуле без служебных кодонов, при Key~0 и при Key~1.
        Для каждой стороны записываются четыре нуклонные величины
        и их остатки по модулю~$37$;
  \item \texttt{prime\_divisibility\_scan.csv} перебирает все
        простые числа до~$97$ для трёх наборов значений:
        групповых величин, одинаковых при Key~0 и Key~1,
        всех величин Key~0 и всех величин Key~1. Для каждого
        простого числа записываются числа делимых экземпляров,
        различных ненулевых значений и различных нечётных ядер,
        а также ориентир $n/p$ для $n$ различных значений. Этот
        ориентир описывает равномерность остатков и не является
        статистической нулевой моделью или $p$-значением;
  \item \texttt{position\_asymmetry\_ncbi.csv} применяет
        классификацию Кето/Амино к каждой из трёх позиций всех
        $27$ таблиц NCBI при Модели~1 и Модели~2.
        Для $T$, $P$, $N$ и
        $\Delta$ записываются направленные разности Кето минус
        Амино, остатки по модулю~$37$ и сводные признаки по
        позициям;
  \item \texttt{representation\_sensitivity.csv} пересчитывает
        $13$ явно перечисленных абсолютных и разностных
        величин при трёх представлениях кодируемой аминокислоты.
        Протонный и нейтронный сдвиги относительно свободной
        нейтральной молекулы равны $(0,0)$ для самой свободной
        молекулы, $(-10,-8)$ для пептидного остатка и $(0,0)$ для
        цвиттер-иона. Для каждой величины записывается результат при
        каждом представлении и его совпадение с базовым значением
        свободной молекулы.
\end{itemize}
Эти четыре файла содержат результаты трёх контролей~---
позиции кодона, выбора модуля и молекулярного представления~---
и межтабличного сравнения разностей Кето/Амино.
Интерпретация приведена в
разделе~\ref{sec:discussion-controls}.

\paragraph{Доступность данных и кода.}
Сопровождающий репозиторий содержит исходники рукописи, пять входных
наборов, $18$ основных вычисляемых наборов, четыре контрольных
набора, две сводные таблицы Монте-Карло
\texttt{mc\_significance.csv} и \texttt{mc\_sensitivity.csv},
шесть Python-скриптов и интерактивный
визуализатор. Полные команды запуска и описание столбцов наборов данных
приведены в README репозитория:
\url{https://github.com/RuslanKhafizov/genetic-code-arithmetic}.
Версии материалов архивируются в Zenodo под концептуальным DOI
\href{https://doi.org/10.5281/zenodo.20360037}
     {10.5281/zenodo.20360037}.

\endgroup

\section{Результаты}
\label{sec:results}

Арифметические свойства рассматриваются здесь на трёх аналитических
уровнях. Первый уровень (раздел~\ref{sec:results-independent}) включает
факты, устанавливаемые без назначения значений служебным кодонам.

Второй уровень (раздел~\ref{sec:results-key0})~--- это наивное
прочтение (Key~0), при котором трём стоп-кодонам приписывается нулевой
вклад, а инициатор $\mathrm{ATG}$ учитывается как метионин.

Раздел~\ref{sec:results-localization} сводит асимметрию ветвей Кето и
Амино к нуклонным разностям пары $\mathrm{Trp}$--$\mathrm{Ile}$ и
рассматривает ограничение чётности на протонные величины.

На третьем уровне раздел~\ref{sec:results-key1-derivation}
выводит Key~1~--- единственное минимальное неотрицательное
целочисленное назначение служебным кодонам, определяемое совместными
условиями баланса нуклонных пулов по оси Кето/Амино и делимости
глобальных пулов на~$37$; эквивалентно, этот баланс дополняет заданную
кодом симметрию $\mathrm{C}/\mathrm{T}$ симметрией
$\mathrm{G}/\mathrm{A}$ в протонных и нейтронных суммах третьей
позиции.

Раздел~\ref{sec:results-key1-consequences} использует это формальное
назначение для распространения арифметического анализа на полную
таблицу из $64$ кодонов и отделяет тождества, непосредственно следующие
из условий вывода, от дополнительных закономерностей, не
использованных при выборе Key~1.

Раздел~\ref{sec:results-synthesis} собирает итоги всех трёх уровней.

Стандартный код остаётся главным объектом анализа на протяжении всей
работы. Альтернативные таблицы трансляции NCBI рассматриваются позднее
для сравнения и оценки специфичности результатов.

\subsection{Результаты, не зависящие от параметризации}
\label{sec:results-independent}

Не зависящие от параметризации результаты для групповых нуклонных
величин в этом подразделе получены при одном из двух условий: либо
область суммирования ограничена пулом без служебных кодонов, либо сама
рассматриваемая группа не содержит таких кодонов.

Эти результаты образуют основу арифметической структуры,
прослеживаемой далее при возвращении служебных кодонов. Нейтронная сумма пула без служебных кодонов кратна~$37$, протонная
меньше ближайшего кратного на единицу. По разности~$\Delta$ это
отклонение прослеживается до кодонов с $\mathrm{G}$ в третьей позиции
внутри Октета~II~--- группы, которая в полной таблице содержит
инициатор $\mathrm{ATG}$ (раздел~\ref{sec:results-sensepool}).
Нуклонная сумма Октета~I кратна~$37$, а все четыре его величины
кратны~$100$
(раздел~\ref{sec:results-octetI}). У пиримидиновых групп полного кода
на решётке шага~$37$ лежит разность~$\Delta$, а отклонения, возникающие
при пересечении с октетами, компенсируются внутри группы
(раздел~\ref{sec:results-pyrimidine}).

\subsubsection{Пул без служебных кодонов}
\label{sec:results-sensepool}

Разбиения $\mathcal{C}_{\mathrm{SE}}$ по осям Кето/Амино и Сильные/Слабые
дают попарно равные нейтронные и протонные суммы вследствие пиримидиновой
синонимии (раздел~\ref{sec:twin-identities}):
\[
\begin{aligned}
N_{\mathrm{SE}}(\mathrm{Keto}) = N_{\mathrm{SE}}(\mathrm{Strong})
  &= 1813 = 49 \cdot 37, \\
N_{\mathrm{SE}}(\mathrm{Amino}) = N_{\mathrm{SE}}(\mathrm{Weak})
  &= 1776 = 48 \cdot 37, \\
P_{\mathrm{SE}}(\mathrm{Keto}) = P_{\mathrm{SE}}(\mathrm{Strong})
  &= 2108 = 57 \cdot 37 - 1, \\
P_{\mathrm{SE}}(\mathrm{Amino}) = P_{\mathrm{SE}}(\mathrm{Weak})
  &= 2072 = 56 \cdot 37.
\end{aligned}
\]
Обе нейтронные суммы лежат на решётке шага~$37$ по отдельности. Протонная
сумма половины Амино попадает на неё точно, а половина Кето не достигает
ближайшего кратного на одну единицу. Складывая половины, получаем для
всего пула
\[
N(\mathcal{C}_{\mathrm{SE}}) = 3589 = 97 \cdot 37,
\qquad
P(\mathcal{C}_{\mathrm{SE}}) = 4180 = 113 \cdot 37 - 1,
\qquad
\Delta(\mathcal{C}_{\mathrm{SE}}) = 591 = 16 \cdot 37 - 1.
\]

Анализ $\Delta$ позволяет локализовать единичное отклонение протонной
суммы пула от ближайшего кратного~$37$ в группе $\{\mathrm{G}\} \cap
\mathrm{Octet~II}$. На осях Кето/Амино и Сильные/Слабые нейтронные суммы
кратны~$37$, поэтому $\Delta = P - N$ имеет тот же остаток, что и
протонная сумма. На оси Пурины/Пиримидины ни $P$, ни $N$ на решётку не
попадают, но их остатки совпадают у Пиримидинов и различаются на единицу у
Пуринов:
\[
\begin{aligned}
\Delta_{\mathrm{SE}}(\mathrm{Keto}) = \Delta_{\mathrm{SE}}(\mathrm{Strong})
  = \Delta_{\mathrm{SE}}(\mathrm{Purine})
  &= 295 = 8 \cdot 37 - 1, \\
\Delta_{\mathrm{SE}}(\mathrm{Amino}) = \Delta_{\mathrm{SE}}(\mathrm{Weak})
  = \Delta(\mathrm{Pyrimidine})
  &= 296 = 8 \cdot 37.
\end{aligned}
\]

Кето, Сильные, Амино и Слабые содержат в пуле по~$30$ кодонов, Пурины~---
$28$, а Пиримидины~--- $32$. Во всех трёх разбиениях при этом получаются
одни и те же значения $\Delta$: $295$ и $296$~--- как при равной мощности
частей $30/30$, так и при неравной мощности $28/32$.

На уровне однонуклеотидных групп ни одна из сумм $P$ и $N$ не кратна $37$,
а единичное отклонение по~$\Delta$ несёт одна группа~--- ветвь $\{\mathrm{G}\}$, общая
для Кето, Сильных и Пуринов:
\[
\Delta_{\mathrm{SE}}(\{\mathrm{G}\}) = 147 = 4 \cdot 37 - 1,
\qquad
\Delta_{\mathrm{SE}}(\{\mathrm{A}\}) =
\Delta(\{\mathrm{C}\}) =
\Delta(\{\mathrm{T}\}) = 148 = 4 \cdot 37.
\]
Кето, Сильные и Пурины складываются из $\{\mathrm{G}\}$ и соответственно
$\{\mathrm{T}\}$, $\{\mathrm{C}\}$ и $\{\mathrm{A}\}$. Поэтому единичное
отклонение, локализованное в $\Delta_{\mathrm{SE}}(\{\mathrm{G}\})$,
прослеживается на всех трёх осях.

В Октете~I все четыре однонуклеотидные группы имеют $\Delta = 75$,
поэтому различие целиком приходится на Октет~II:
\[
\Delta_{\mathrm{SE}}(\{\mathrm{G}\} \cap \mathrm{Octet~II}) = 72,
\qquad
\Delta_{\mathrm{SE}}(\{\mathrm{A}\} \cap \mathrm{Octet~II}) =
\Delta(\{\mathrm{C}\} \cap \mathrm{Octet~II}) =
\Delta(\{\mathrm{T}\} \cap \mathrm{Octet~II}) = 73.
\]
В пуле группы $\{\mathrm{G}\} \cap \mathrm{Octet~II}$ и $\{\mathrm{A}\}
\cap \mathrm{Octet~II}$ содержат по шесть кодонов, а группы
$\{\mathrm{C}\} \cap \mathrm{Octet~II}$ и $\{\mathrm{T}\} \cap
\mathrm{Octet~II}$~--- по восемь. Наблюдаемое распределение значений $72$
и $73$ тем самым не сводится к мощности групп.

При переходе к полной 64-кодонной таблице эти две шестикодонные ветви
дополняются разными парами служебных кодонов: $\{\mathrm{A}\} \cap
\mathrm{Octet~II}$~--- стоп-кодонами $\mathrm{TAA}$ и $\mathrm{TGA}$, а
$\{\mathrm{G}\} \cap \mathrm{Octet~II}$~--- стоп-кодоном $\mathrm{TAG}$ и
инициатором $\mathrm{ATG}$. Как будет получено при выводе Key~1
(раздел~\ref{sec:results-key1-derivation}), формальный $\Delta$-вклад
инициатора $\mathrm{ATG}$ при этом назначении равен~$1$, а каждого
стоп-кодона~--- $0$. Поэтому при Key~1
\[
\Delta_{1}(\{\mathrm{G}\} \cap \mathrm{Octet~II})
  = 72 + 1 + 0 = 73,
\qquad
\Delta_{1}(\{\mathrm{A}\} \cap \mathrm{Octet~II})
  = 73 + 0 + 0 = 73.
\]
Недостающая единица возвращается тем самым именно через инициатор
$\mathrm{ATG}$, выравнивая обе ветви. Поскольку $\mathrm{ATG}$ входит в
ветвь $\{\mathrm{G}\}$, Кето, Сильные и Пурины, тот же единичный вклад
одновременно компенсирует отклонение во всех этих группах и в значении
$\Delta$ для полной таблицы при Key~1.

Полная иерархия групп для пула без служебных кодонов вместе с
соответствующими нуклонными величинами показана на Рис.~\ref{fig:pool60}.

\begin{figure}[p]
\centering
\includegraphics[width=0.72\textwidth]{pool60_structure.png}
\caption{Структура 60-кодонного пула без служебных кодонов
$\mathcal{C}_{\mathrm{SE}}$ (\texttt{pool60}). Для каждой группы все четыре
величины вычисляются только по входящим в неё неслужебным кодонам; три
стоп-кодона и инициатор $\mathrm{ATG}$ исключены, поэтому каждая
показанная величина не зависит от параметризации. Приведены нуклонная
сумма~$T$, протонная
сумма~$P$, нейтронная сумма~$N$ и число
атомов водорода $\Delta = P - N$. Оформление ячейки: название группы в
левом верхнем углу, число оставшихся в ней кодонов в правом верхнем,
идентификатор из \texttt{codon\_groups.csv} в правом нижнем; пары
групп-близнецов с совпадающими значениями сведены в одну ячейку, второе
название стоит в левом нижнем углу, а идентификаторы обеих групп
перечислены через запятую в правом нижнем. Зелёная, синяя и фиолетовая
подсветка обозначают делимость соответственно на~$37$, $111 = 3 \cdot 37$
и $999 = 27 \cdot 37$; отмечается наибольший из применимых модулей.
Рис.~\ref{fig:groups-key0} и~\ref{fig:groups-key1} построены по той же
сетке и с тем же цветовым кодом, но для полной таблицы из $64$ кодонов.
Все три рисунка сгенерированы сопровождающим \texttt{visualization.html}.}
\label{fig:pool60}
\end{figure}

\subsubsection{Арифметика Октета~I}
\label{sec:results-octetI}

Октет~I (id~12) не содержит служебных кодонов, поэтому его нуклонные
величины не зависят от параметризации:
\[
\begin{aligned}
T(\mathrm{Octet~I}) &= 3700 = 100 \cdot 37, &
P(\mathrm{Octet~I}) &= 2000 = 2 \cdot 1000, \\
N(\mathrm{Octet~I}) &= 1700 = 17 \cdot 100, &
\Delta(\mathrm{Octet~I}) &= 300 = 3 \cdot 100.
\end{aligned}
\]
Все четыре величины кратны~$100$, а протонная сумма
$P(\mathrm{Octet~I})$ также кратна~$1000$. Кроме того, нуклонная
сумма $T(\mathrm{Octet~I})$ кратна~$37$.

Октет~I состоит из восьми полностью четырёхкратно вырожденных кодонных
семейств. Каждая из шести его $16$-кодонных групп (id~13--18) содержит
по два кодона каждого семейства, поэтому каждая её нуклонная величина
равна половине соответствующей величины всего октета:
\[
T = 1850 = 50 \cdot 37, \qquad
P = 1000 = 10^{3}, \qquad
N = 850, \qquad
\Delta = 150.
\]
Каждая из четырёх $8$-кодонных групп (id~19--22) содержит по одному
кодону каждого семейства, поэтому её величины составляют четверть
значений всего октета:
\[
T = 925 = 25 \cdot 37, \qquad
P = 500 = 5 \cdot 100, \qquad
N = 425, \qquad
\Delta = 75 = 2 \cdot 37 + 1.
\]
Таким образом, точное масштабирование сохраняет на всех трёх уровнях
иерархии Октета~I делимость нуклонной суммы $T$ на~$37$ и единый
рациональный профиль с общим знаменателем~$37$:
\[
\frac{P}{T} = \frac{20}{37},
\qquad
\frac{N}{T} = \frac{17}{37},
\qquad
\frac{\Delta}{T} = \frac{3}{37}.
\]

Протонная и нейтронная суммы всего октета имеют остатки $+2$ и $-2$
по модулю~$37$, которые взаимно сокращаются в $T = P + N$. Нейтронный
остаток $-2$ Октета~I, в свою очередь, компенсируется остатком $+2$
нейтронной суммы Октета~II в пуле без служебных кодонов. Сложение
нейтронных сумм двух октетов даёт нейтронную сумму пула
$N(\mathcal{C}_{\mathrm{SE}})$, полученную в
разделе~\ref{sec:results-sensepool}. Оба сокращения представлены
следующими равенствами:
\[
\begin{aligned}
P(\mathrm{Octet~I})
  &= 2000 = 54 \cdot 37 + 2, \\
N(\mathrm{Octet~I})
  &= 1700 = 46 \cdot 37 - 2, \\
N_{\mathrm{SE}}(\mathrm{Octet~II})
  &= 1889 = 51 \cdot 37 + 2, \\
N(\mathcal{C}_{\mathrm{SE}})
  &= N(\mathrm{Octet~I})
   + N_{\mathrm{SE}}(\mathrm{Octet~II}) \\
  &= 1700 + 1889 = 3589 = 97 \cdot 37.
\end{aligned}
\]

На $16$-кодонном уровне точное деление пополам даёт остатки $+1$ для
$P$ и $-1$ для $N$ по модулю~$37$.
Из четырёх $8$-кодонных групп две пиримидиновые~---
$\{\mathrm{C}\} \cap \mathrm{Octet~I}$ (id~19) и
$\{\mathrm{T}\} \cap \mathrm{Octet~I}$ (id~21)~--- образуют группу
Пиримидины в Октете~I (id~18), которая далее входит в пиримидиновый
каскад раздела~\ref{sec:results-pyrimidine}.

\subsubsection{Арифметика групп, ограниченных пиримидинами}
\label{sec:results-pyrimidine}

В силу пиримидиновой синонимии
(раздел~\ref{sec:twin-identities}) группы $\{\mathrm{C}\}$ и
$\{\mathrm{T}\}$ имеют одинаковые нуклонные величины на уровне
полного кода и после пересечения с каждым октетом. Вместе они образуют
группу Пиримидины (id~7). Эти группы и их октетные пересечения не
содержат служебных кодонов, поэтому их значения не зависят от
параметризации. В отличие от иерархии Октета~I, полная группа
Пиримидины охватывает оба октета и включает частично вырожденные
кодонные семейства. Поэтому не все значения её нижних уровней
получаются единым точным масштабированием; ниже они рассматриваются
раздельно.

Для полной группы Пиримидины нуклонные величины равны
\[
\begin{aligned}
T(\mathrm{Pyrimidine}) &= 4096 = 2^{12}, &
P(\mathrm{Pyrimidine}) &= 2196, \\
N(\mathrm{Pyrimidine}) &= 1900 = 19 \cdot 100, &
\Delta(\mathrm{Pyrimidine}) &= 296 = 8 \cdot 37.
\end{aligned}
\]
Здесь $\Delta$ удовлетворяет основному заранее выбранному критерию
делимости на~$37$, а $N = 1900$~--- критерию десятичной круглости.
Точная степень двойки $T = 2^{12}$ является дополнительным описательным
наблюдением.

По тождеству C/T две $16$-кодонные группы $\{\mathrm{C}\}$ (id~8) и
$\{\mathrm{T}\}$ (id~10) имеют одинаковые нуклонные величины.
Поскольку они образуют разбиение группы Пиримидины, каждое значение
равно половине соответствующего значения полной группы:
\[
\begin{aligned}
T(\{\mathrm{C}\}) = T(\{\mathrm{T}\})
  &= 2048 = 2^{11}, &
P(\{\mathrm{C}\}) = P(\{\mathrm{T}\})
  &= 1098, \\
N(\{\mathrm{C}\}) = N(\{\mathrm{T}\})
  &= 950, &
\Delta(\{\mathrm{C}\}) = \Delta(\{\mathrm{T}\})
  &= 148 = 4 \cdot 37.
\end{aligned}
\]
Весь профиль этого уровня тем самым получается точным масштабированием
и не образует отдельного набора независимых закономерностей.

На уровне $8$ кодонов все четыре группы Октета~I, в том числе
$\{\mathrm{C}\} \cap \mathrm{Octet~I}$ (id~19) и
$\{\mathrm{T}\} \cap \mathrm{Octet~I}$ (id~21), имеют
$\Delta = 75$ (раздел~\ref{sec:results-octetI}). Соответствующие группы
$\{\mathrm{C}\} \cap \mathrm{Octet~II}$ (id~30) и
$\{\mathrm{T}\} \cap \mathrm{Octet~II}$ (id~32) имеют
$\Delta = 73$. Их объединения внутри октетов образуют группы
$\mathrm{Pyrimidine} \cap \mathrm{Octet~I}$ (id~18) и
$\mathrm{Pyrimidine} \cap \mathrm{Octet~II}$ (id~29). Эти значения
связаны следующим каскадом:
\[
\begin{aligned}
\Delta(\{\mathrm{C}\} \cap \mathrm{Octet~I})
  = \Delta(\{\mathrm{T}\} \cap \mathrm{Octet~I})
  &= 75 = 2 \cdot 37 + 1, \\
\Delta(\{\mathrm{C}\} \cap \mathrm{Octet~II})
  = \Delta(\{\mathrm{T}\} \cap \mathrm{Octet~II})
  &= 73 = 2 \cdot 37 - 1, \\
\Delta(\mathrm{Pyrimidine} \cap \mathrm{Octet~I})
  &= 2 \cdot 75 = 150 = 4 \cdot 37 + 2, \\
\Delta(\mathrm{Pyrimidine} \cap \mathrm{Octet~II})
  &= 2 \cdot 73 = 146 = 4 \cdot 37 - 2, \\
\Delta(\mathrm{Pyrimidine})
  &= 150 + 146 = 296 = 8 \cdot 37.
\end{aligned}
\]
Единичные отклонения противоположных знаков, $+1$ и $-1$, при
сложении групп внутри каждого октета удваиваются до $+2$ и $-2$, а при
сложении октетов взаимно компенсируются. Поэтому значения $150$, $146$
и $296$ связаны одной аддитивной схемой и не являются независимыми
результатами.

Связь пиримидинового каскада с паттерном единичных отклонений
протонных и нейтронных сумм рассматривается при Key~0
(раздел~\ref{sec:results-key0-delta}).

Пиримидиновый каскад следует также отличать от единичного отклонения
$\Delta_{\mathrm{SE}}(\{\mathrm{G}\})$ от решётки шага~$37$
(раздел~\ref{sec:results-sensepool}). Пиримидиновые отклонения
компенсируются между октетами, тогда как отклонение группы
$\{\mathrm{G}\}$ не компенсируется при объединении с другой
однонуклеотидной группой и передаётся группам Кето, Сильные и Пурины.
Поскольку нейтронные суммы групп Кето и Сильные в пуле без служебных
кодонов кратны~$37$, из тождества $P=N+\Delta$ следует, что их
протонные суммы имеют то же единичное отклонение $-1$.

\subsection{Результаты при Key~0}
\label{sec:results-key0}

При наивном прочтении (Key~0) суммирование ведётся по полной
$64$-кодонной таблице: вклад стоп-кодонов равен нулю, а инициатор
$\mathrm{ATG}$ учитывается как метионин
(раздел~\ref{sec:parametrization}). Поэтому при переходе от пула без
служебных кодонов единственный ненулевой вклад вносит $\mathrm{ATG}$:
\[
\bigl(T(\mathrm{Met}),P(\mathrm{Met}),N(\mathrm{Met}),
\Delta(\mathrm{Met})\bigr)=(149,80,69,11).
\]

Отклонения от ближайших кратных~$37$ у пула без служебных кодонов,
у метионина и у полной $64$-кодонной таблицы связаны как
\[
(-1,-1,0,-1)+(+1,+6,-5,+11)=(0,+5,-5,+10).
\]
Вклад $\mathrm{ATG}$ выводит нуклонную сумму $T$ точно на решётку
кратных~$37$: $149=4\cdot37+1$ точно компенсирует единичное
отклонение $-1$, но одновременно уводит с этой решётки нейтронную
сумму, которая в пуле кратна~$37$
(раздел~\ref{sec:results-sensepool}).

Ниже показано, как это преобразование распространяется по вложенным
группам, содержащим $\mathrm{ATG}$. В группах Амино и Слабые все
возвращённые позиции заняты стоп-кодонами, поэтому суммы полных
$32$-кодонных групп численно совпадают с суммами по $30$ кодонам
тех же ветвей в пуле без служебных кодонов. Эти суммы уже лежат на
решётке шага~$37$ (раздел~\ref{sec:results-sensepool}), так что
делимость при переходе к полной таблице сохраняется. При Key~0
возвращение служебных позиций делает равными по размеру группы,
образующие две стороны зеркального паттерна единичных отклонений.
Паттерн проявляется на двух уровнях: на уровне $16$ кодонов~--- в
протонных и нейтронных суммах, а на уровне $8$ кодонов~--- в
разности $\Delta$. В Октете~II на уровне $8$ кодонов этот паттерн, ранее установленный
для пиримидиновых групп
(раздел~\ref{sec:results-pyrimidine}), охватывает и группу
$\{\mathrm{A}\}$: её значение $\Delta$ уже совпадало с ними в пуле,
а возвращение служебных позиций делает группы равными и по размеру. Единственным
исключением остаётся ветвь $\{\mathrm{G}\}$, содержащая
$\mathrm{ATG}$.

Общая картина групповых величин при Key~0 показана на
Рис.~\ref{fig:groups-key0}.

\begin{figure}[!htbp]
  \centering
  \includegraphics[width=0.78\textwidth]{groups_key0.png}
  \caption{Нуклонные величины $T$, $P$, $N$ и $\Delta$ всех групп
  кодонов при Key~0. Сетка, свёртка пар групп-близнецов и цветовой
  код делимости те же, что на Рис.~\ref{fig:pool60}, но величины
  вычислены по полной таблице из $64$ кодонов. Сгенерировано
  сопровождающим \texttt{visualization.html}.}
  \label{fig:groups-key0}
\end{figure}

\subsubsection{Делимость \texorpdfstring{$T$}{T} на
\texorpdfstring{$37$}{37} на верхних уровнях разбиения}

Вклад $\mathrm{ATG}$ приходится на Октет~II. При Key~0 его четыре
$8$-кодоновые группы имеют следующий профиль:
\[
\begin{aligned}
T(\{\mathrm{C}\}\cap\mathrm{Octet~II})
  =T(\{\mathrm{T}\}\cap\mathrm{Octet~II})
  &=1123=30\cdot37+13, \\
T_{0}(\{\mathrm{A}\}\cap\mathrm{Octet~II})
  &=875=24\cdot37-13, \\
T_{0}(\{\mathrm{G}\}\cap\mathrm{Octet~II})
  &=1097=30\cdot37-13.
\end{aligned}
\]
Пиримидиновые ветви отклоняются от ближайшего кратного на $+13$,
пуриновые~--- на $-13$. В пуле без служебных кодонов отклонение
ветви $\{\mathrm{G}\}$ отличается от отклонения ветви
$\{\mathrm{A}\}$ на единицу:
\[
T_{\mathrm{SE}}(\{\mathrm{G}\}\cap\mathrm{Octet~II})
  =948=26\cdot37-14.
\]
Вклад $\mathrm{ATG}$ меняет отклонение ветви $\{\mathrm{G}\}$ с
$-14$ до $-13$, так что обе пуриновые ветви имеют одинаковое
отклонение.

Среди групп Октета~II нуклонная сумма делится на~$37$ в точности у
тех, в которых поровну пуриновых и пиримидиновых ветвей:
\[
\begin{aligned}
T_{0}(\mathrm{Keto}\cap\mathrm{Octet~II})
  =T_{0}(\mathrm{Strong}\cap\mathrm{Octet~II})
  &=2220=60\cdot37, \\
T_{0}(\mathrm{Amino}\cap\mathrm{Octet~II})
  =T_{0}(\mathrm{Weak}\cap\mathrm{Octet~II})
  &=1998=54\cdot37, \\
T_{0}(\mathrm{Octet~II})
  &=4218=114\cdot37.
\end{aligned}
\]
Пересечения групп Пурины и Пиримидины с Октетом~II состоят из
ветвей одного типа, поэтому отклонения ветвей не компенсируются, а
удваиваются относительно сумм соответствующих кратных~$37$.

Каждая $16$-кодоновая половина химической оси в Октете~I имеет
$T=1850=50\cdot37$ (раздел~\ref{sec:results-octetI}). Прибавление
соответствующей половины переносит делимость на полные группы осей
Кето/Амино и Сильные/Слабые, а сложение двух октетов~--- на всю
таблицу:
\[
\begin{aligned}
T_{0}(\mathrm{Keto})=T_{0}(\mathrm{Strong})
  &=1850+2220=4070=110\cdot37, \\
T_{0}(\mathrm{Amino})=T_{0}(\mathrm{Weak})
  &=1850+1998=3848=104\cdot37, \\
T_{0}(\mathrm{All})
  &=3700+4218=7918=214\cdot37.
\end{aligned}
\]
Перечисленные делимости следуют из зеркальных отклонений $+13/-13$
между ветвями Октета~II и из делимости Октета~I, а не образуют
независимого набора результатов.

Делимости $T(\mathrm{Octet~I})=3700$,
$T_{0}(\mathrm{Octet~II})=4218$ и
$T_{0}(\mathrm{Keto}\cap\mathrm{Octet~II})
=T_{0}(\mathrm{Strong}\cap\mathrm{Octet~II})=2220$
соответствуют в принятом здесь групповом формализме результатам
Щербака и Макукова~\cite{makukov2013} и Панова и
Филатова~\cite{panov2024}; подробнее см.
раздел~\ref{sec:discussion-prior-work}. Полные значения групповых
величин при Key~0 приведены в
\texttt{key0\_nucleon-data.csv}.

\subsubsection{Выравнивание групп Амино и Слабые с решёткой
шага~\texorpdfstring{$37$}{37}}
\label{sec:results-key0-aminoweak}

Среди шести $32$-кодоновых групп трёх химических осей только Амино
(id~4) и Слабые (id~5) имеют все четыре нуклонные величины,
кратные~$37$:
\[
\begin{aligned}
T_{0}(\mathrm{Amino})=T_{0}(\mathrm{Weak})
  &=3848=104\cdot37, \\
P_{0}(\mathrm{Amino})=P_{0}(\mathrm{Weak})
  &=2072=56\cdot37, \\
N_{0}(\mathrm{Amino})=N_{0}(\mathrm{Weak})
  &=1776=48\cdot37, \\
\Delta_{0}(\mathrm{Amino})=\Delta_{0}(\mathrm{Weak})
  &=296=8\cdot37.
\end{aligned}
\]
Делимость $T$ получена выше. В Октете~I половины обеих групп имеют
остатки $+1$ для $P$ и $-1$ для $N$
(раздел~\ref{sec:results-octetI}), в Октете~II~--- противоположные:
\[
\begin{aligned}
P(\mathrm{Amino}\cap\mathrm{Octet~I})
  =P(\mathrm{Weak}\cap\mathrm{Octet~I})
  &=1000=27\cdot37+1, \\
P_{0}(\mathrm{Amino}\cap\mathrm{Octet~II})
  =P_{0}(\mathrm{Weak}\cap\mathrm{Octet~II})
  &=1072=29\cdot37-1, \\
N(\mathrm{Amino}\cap\mathrm{Octet~I})
  =N(\mathrm{Weak}\cap\mathrm{Octet~I})
  &=850=23\cdot37-1, \\
N_{0}(\mathrm{Amino}\cap\mathrm{Octet~II})
  =N_{0}(\mathrm{Weak}\cap\mathrm{Octet~II})
  &=926=25\cdot37+1.
\end{aligned}
\]
При сложении половин эти остатки взаимно компенсируются, а вместе с
ними и остаток разности $\Delta=P-N$.

Для каждой из групп $G\in\{\mathrm{Amino},\mathrm{Weak}\}$
отношения к нуклонной сумме при Key~0 образуют профиль
с общим знаменателем~$13$:
\[
\frac{P_0(G)}{T_0(G)}=\frac{7}{13},\qquad
\frac{N_0(G)}{T_0(G)}=\frac{6}{13},\qquad
\frac{\Delta_0(G)}{T_0(G)}=\frac{1}{13}.
\]

\subsubsection{Единичные отклонения по
\texorpdfstring{$\Delta$}{Delta} и ветвь
\texorpdfstring{$\{\mathrm{G}\}$}{G}}
\label{sec:results-key0-delta}

В Октете~I все четыре ветви имеют
$\Delta=75=2\cdot37+1$
(раздел~\ref{sec:results-octetI}). В Октете~II отклонение
противоположно, и при Key~0 его несут три ветви из четырёх:
\[
\begin{aligned}
\Delta(\{\mathrm{C}\}\cap\mathrm{Octet~II})
  =\Delta(\{\mathrm{T}\}\cap\mathrm{Octet~II})
  =\Delta_{0}(\{\mathrm{A}\}\cap\mathrm{Octet~II})
  &=73=2\cdot37-1, \\
\Delta_{0}(\{\mathrm{G}\}\cap\mathrm{Octet~II})
  &=83=2\cdot37+9.
\end{aligned}
\]
Равенство ветвей $\{\mathrm{C}\}$ и $\{\mathrm{T}\}$ следует из
пиримидиновой синонимии
(раздел~\ref{sec:results-pyrimidine}), тогда как совпадение с
пуриновой ветвью $\{\mathrm{A}\}$ является дополнительным числовым
фактом. При Key~0 её две возвращённые позиции заняты стоп-кодонами,
поэтому значение пула $\Delta=73$ сохраняется, а размер группы
увеличивается до восьми кодонов. Ветвь $\{\mathrm{G}\}$ отличается
уже в пуле без служебных кодонов, где имеет $\Delta=72$ против $73$ у остальных ветвей;
вклад $\mathrm{ATG}$ изменяет это значение до $83$
(раздел~\ref{sec:results-sensepool}).

Прибавление ветвей Октета~I даёт значения однонуклеотидных групп
полной таблицы:
\[
\Delta(\{\mathrm{C}\})=\Delta(\{\mathrm{T}\})
  =\Delta_{0}(\{\mathrm{A}\})=148=4\cdot37,
\qquad
\Delta_{0}(\{\mathrm{G}\})=158=4\cdot37+10.
\]
Ветвь $\{\mathrm{G}\}$ входит в Кето, Сильные и Пурины, тогда как
остальные три однонуклеотидные группы лежат на решётке. Поэтому
остаток $+10$ передаётся всем трём осям и полной таблице:
\[
\begin{aligned}
\Delta_{0}(\mathrm{Keto})
  =\Delta_{0}(\mathrm{Strong})
  =\Delta_{0}(\mathrm{Purine})
  &=306=8\cdot37+10, \\
\Delta_{0}(\mathrm{All})
  &=602=16\cdot37+10.
\end{aligned}
\]
Тем самым отклонение полной таблицы по $\Delta$, приведённое в
начале раздела, локализовано в ветви $\{\mathrm{G}\}$.

Из $8$-кодонового паттерна следуют и остатки $P$ и $N$ на уровне
$16$ кодонов (раздел~\ref{sec:results-key0-aminoweak}). В половинах
групп Амино и Слабые два остатка $\Delta\equiv+1$ в Октете~I
складываются в $+2$, а два остатка $\Delta\equiv-1$ в
Октете~II~--- в $-2$. Поскольку соответствующие суммы $T$
кратны~$37$, из $2P=T+\Delta$ и $2N=T-\Delta$ следуют
$(P,N)\equiv(+1,-1)\pmod{37}$ в Октете~I и
$(P,N)\equiv(-1,+1)\pmod{37}$ в Октете~II.

\subsection{Локализация асимметрии Кето/Амино и протонная чётность}
\label{sec:results-localization}

Асимметрия ветвей Кето и Амино видна на обоих рассмотренных уровнях.
Ниже она показана сначала на полной таблице при Key~0, где сводится к
нуклонным величинам трёх молекул, затем в пуле без служебных кодонов,
где остаются две.

\subsubsection{Локализация асимметрии Кето/Амино}
\label{sec:results-key1-trpile}

На полной таблице при Key~0 разности нуклонных величин ветвей Кето и
Амино равны
\[
\begin{aligned}
T_{0}(\mathrm{Keto}) - T_{0}(\mathrm{Amino})
  &= 4070 - 3848 = 222 = 6 \cdot 37, \\
P_{0}(\mathrm{Keto}) - P_{0}(\mathrm{Amino})
  &= 2188 - 2072 = 116, \\
N_{0}(\mathrm{Keto}) - N_{0}(\mathrm{Amino})
  &= 1882 - 1776 = 106, \\
\Delta_{0}(\mathrm{Keto}) - \Delta_{0}(\mathrm{Amino})
  &= 306 - 296 = 10.
\end{aligned}
\]

\paragraph{Сведение на уровне семейств.}
В силу пиримидиновой синонимии
(раздел~\ref{sec:twin-identities}) во всех $16$ кодонных семействах
кодоны $\mathrm{NNT}$ и $\mathrm{NNC}$ кодируют одну и ту же
аминокислоту. В $13$ семействах, не содержащих служебных кодонов,
совпадают также продукты $\mathrm{NNG}$ и $\mathrm{NNA}$. Поэтому
ветвь Кето, представленная парой $\mathrm{NNG}$ и $\mathrm{NNT}$,
и ветвь Амино, представленная парой $\mathrm{NNA}$ и
$\mathrm{NNC}$, вносят равный нуклонный вклад независимо от того,
какие именно аминокислоты кодирует данное семейство.

Семейство $\mathrm{TA}$ также вносит равный вклад: его
пиримидиновая пара $\mathrm{TAC}$ и $\mathrm{TAT}$ состоит из
двух кодонов $\mathrm{Tyr}$, а пуриновая пара $\mathrm{TAG}$ и
$\mathrm{TAA}$~--- из двух стоп-кодонов, по одному в каждой
ветви. Последние сокращаются при одинаковых назначениях
этим двум стоп-кодонам.

Асимметрия остаётся в двух семействах, и в каждом из них одна пара
кодонов сокращается, а вторая нет. В $\mathrm{AT}$ сокращаются
$\mathrm{ATT}$ и $\mathrm{ATC}$,
кодирующие $\mathrm{Ile}$ в ветвях Кето и Амино; остаётся пара
$\mathrm{ATG}/\mathrm{ATA}$: инициатор против третьего кодона
$\mathrm{Ile}$. В $\mathrm{TG}$ сокращаются $\mathrm{TGT}$ и
$\mathrm{TGC}$, кодирующие $\mathrm{Cys}$; остаётся пара
$\mathrm{TGG}/\mathrm{TGA}$: кодон $\mathrm{Trp}$ против
стоп-кодона.

При Key~0 инициатор учитывается как кодируемый им метионин, а
стоп-кодон $\mathrm{TGA}$ вклада не даёт. Поэтому из двух оставшихся
пар для $X\in\{T,P,N,\Delta\}$
\[
X_{0}(\mathrm{Keto})-X_{0}(\mathrm{Amino})
=
X(\mathrm{Trp})+X(\mathrm{Met})-X(\mathrm{Ile}),
\]
то есть асимметрия полной таблицы сводится к нуклонным величинам трёх
молекул; в частности,
\[
T_{0}(\mathrm{Keto}) - T_{0}(\mathrm{Amino})
  = 204 + 149 - 131 = 222 = 6 \cdot 37.
\]

\paragraph{Сведение к паре $\mathrm{Trp}$--$\mathrm{Ile}$.}
Инициатор $\mathrm{ATG}$~--- единственный служебный кодон, дающий при
Key~0 вклад в эти разности. В пуле без служебных кодонов он исключён
из сумм вместе с тремя стоп-кодонами, и асимметрия сводится
к нуклонным разностям одной пары молекул,
$\mathrm{Trp}$--$\mathrm{Ile}$:
\[
\begin{aligned}
T_{\mathrm{SE}}(\mathrm{Keto}) - T_{\mathrm{SE}}(\mathrm{Amino})
  &= 3921 - 3848 = 73, \\
P_{\mathrm{SE}}(\mathrm{Keto}) - P_{\mathrm{SE}}(\mathrm{Amino})
  &= 2108 - 2072 = 36, \\
N_{\mathrm{SE}}(\mathrm{Keto}) - N_{\mathrm{SE}}(\mathrm{Amino})
  &= 1813 - 1776 = 37, \\
\Delta_{\mathrm{SE}}(\mathrm{Keto})
  - \Delta_{\mathrm{SE}}(\mathrm{Amino})
  &= 295 - 296 = -1.
\end{aligned}
\]
То есть для $X\in\{T,P,N,\Delta\}$
\[
X_{\mathrm{SE}}(\mathrm{Keto})
-
X_{\mathrm{SE}}(\mathrm{Amino})
=
X(\mathrm{Trp})-X(\mathrm{Ile}).
\]
Это сведение определяется структурой синонимии и расположением
служебных кодонов.

По тождествам групп-близнецов (раздел~\ref{sec:twin-identities})
$Q_{\mathrm{SE}}(\mathrm{Strong}) = Q_{\mathrm{SE}}(\mathrm{Keto})$ и
$Q_{\mathrm{SE}}(\mathrm{Weak}) = Q_{\mathrm{SE}}(\mathrm{Amino})$ для
$Q\in\{T,P,N,\Delta\}$. Поэтому асимметрия оси Сильные/Слабые
описывается теми же разностями, что и Кето/Амино, и также сводится к
паре $\mathrm{Trp}$--$\mathrm{Ile}$.

\paragraph{Локальный и глобальный масштабы.}
Для локализованной выше пары соответствующие разности непосредственно
следуют из нуклонных чисел аминокислот:
\[
N(\mathrm{Trp})-N(\mathrm{Ile})=96-59=37,
\qquad
P(\mathrm{Trp})-P(\mathrm{Ile})=108-72=36.
\]
Протонная и нейтронная разности отличаются на единицу, что
соответствует одному атому водорода: триптофан содержит $12$ атомов
водорода, а изолейцин~--- $13$.

Те же остатки по модулю~$37$ были ранее установлены на глобальном
уровне для всего пула без служебных кодонов
(раздел~\ref{sec:results-sensepool}):
\[
\begin{array}{c@{\qquad}cc}
 & \text{глобально} & \text{локально} \\
N &  0 &  0 \\
P & -1 & -1
\end{array}
\]

Глобальные и локальные факты не следуют друг из друга: первые
относятся ко всему пулу без служебных кодонов, вторые~--- к
нуклонному профилю пары $\mathrm{Trp}$--$\mathrm{Ile}$, к которой
сводится асимметрия Кето/Амино в этом пуле.

В дальнейшем локальные разности входят в условия баланса служебных
кодонов, а глобальные остатки~--- в условия делимости при
аналитическом выводе Key~1
(раздел~\ref{sec:results-key1-derivation}).

\paragraph{Арифметическая редкость и структурная единственность пары
Trp--Ile.}
Нейтронная и протонная разности этой пары редки в фиксированном наборе
канонических аминокислот.

\emph{Значение $|\delta N| = 37$.} Среди $190$ пар канонических
аминокислот значение $|\delta N| = 37$ реализуется лишь для двух
пар: $\mathrm{Trp}$--$\mathrm{Ile}$ и
$\mathrm{Trp}$--$\mathrm{Leu}$. Последняя пара даёт то же
значение, поскольку $\mathrm{Leu}$ и $\mathrm{Ile}$ являются
изомерами с тождественными нуклонными числами $P = 72$, $N = 59$;
пара, структурно реализуемая через семейства $\mathrm{AT}$ и
$\mathrm{TG}$,~--- это $\mathrm{Trp}$--$\mathrm{Ile}$. Полное
распределение $|\delta N|$ по $190$ парам приведено в
\texttt{amino\_acids\_neutron\_differences.csv}.

\emph{Значение $|\delta P| = 36$.} Среди $190$ пар канонических
аминокислот значение $|\delta P| = 36$ реализуется лишь для тех
же двух пар~--- $\mathrm{Trp}$--$\mathrm{Ile}$ и
$\mathrm{Trp}$--$\mathrm{Leu}$. Полное распределение
$|\delta P|$ по $190$ парам приведено в
\texttt{amino\_acids\_proton\_differences.csv}.

\emph{Значение $|\delta T| = 73$.} Для пары
$\mathrm{Trp}$--$\mathrm{Ile}$ это значение следует из уже
приведённых протонной и нейтронной разностей, имеющих один знак.
При этом среди всех $190$ пар канонических аминокислот
$|\delta T| = 73$ реализуется только для
$\mathrm{Trp}$--$\mathrm{Ile}$ и
$\mathrm{Trp}$--$\mathrm{Leu}$. Полное распределение
$|\delta T|$ приведено в
\texttt{amino\_acids\_nucleon\_differences.csv}.

\subsubsection{Протонная чётность}
\label{sec:results-proton-parity}

Каждая из $20$ канонических аминокислот имеет чётное суммарное число
протонов: все значения $P$ в
\texttt{amino\_acids\_nucleons.csv} чётны. В основном электронном
состоянии нейтральные молекулы канонических аминокислот являются
замкнутооболочечными системами, в которых электроны занимают
орбитали попарно; поэтому их общее число чётно. Поскольку для
нейтральной молекулы число электронов равно суммарному числу протонов
её ядер, отсюда следует чётность~$P$.

Следовательно, чётны и любая протонная сумма по аминокислотам, и любая
попарная разность, и ни одна такая величина не может равняться
нечётному кратному~$37$. Так, $|\delta P(\mathrm{Trp},\mathrm{Ile})|$
не может равняться $37$: ближайшие к $37$ чётные значения~--- $36$
и~$38$, а фактическая разность для локализованной пары равна $36$.
По той же причине протонная сумма пула без служебных кодонов не может
равняться ближайшему сверху кратному $113 \cdot 37$, которое нечётно
(раздел~\ref{sec:results-sensepool}).

Чётность задаёт только запрет. Она не определяет ни направление, ни
величину отклонения: то, что и локальная разность, и глобальная сумма
оказываются на единицу ниже соответствующего кратного~$37$, следует
из конкретных нуклонных чисел.

\subsection{Аналитический вывод Key~1}
\label{sec:results-key1-derivation}

Результаты разделов~\ref{sec:results-independent}
и~\ref{sec:results-key0} показывают, что проявление
арифметических закономерностей связано с тем, какие
служебные кодоны содержит группа.

Октет~I и группы, ограниченные пиримидинами, не содержат
служебных кодонов и при этом демонстрируют арифметическую
упорядоченность (разделы~\ref{sec:results-octetI}
и~\ref{sec:results-pyrimidine}). Ветви Амино и Слабые
содержат стоп-кодоны, но не инициатор $\mathrm{ATG}$: все
четыре их нуклонные величины кратны~$37$ уже в пуле без
служебных кодонов, а нулевой вклад стоп-кодонов их не
меняет, поэтому при Key~0 делимость сохраняется
(раздел~\ref{sec:results-key0-aminoweak}).
В ветвях Кето и Сильные, содержащих инициатор
$\mathrm{ATG}$, при Key~0 протонные и нейтронные суммы, а
также разность $\Delta$ отклоняются от решётки с
шагом~$37$. Отклонение по $\Delta$ локализовано в ветви
$\{\mathrm{G}\}$, содержащей $\mathrm{ATG}$: остальные три
однонуклеотидные группы имеют одинаковое значение
$\Delta$, кратное~$37$
(раздел~\ref{sec:results-key0-delta}).

Это указывает на возможность иной параметризации
служебных кодонов, согласующей арифметические
закономерности рассмотренных групп. В настоящем разделе
мы аналитически выводим такую параметризацию из условий
симметрии, делимости и минимальности. Полученное
назначение служит формальным средством согласования сумм
полной таблицы; исходные арифметические закономерности
установлены в пуле без служебных кодонов. Специфичность
результата среди природных кодов рассматривается в
разделе~\ref{sec:results-key1-ncbi}.

Для искомого назначения $K$ обозначим через
$(p_{\text{stop}}, n_{\text{stop}})$ общие параметры трёх
стоп-кодонов, а через
$(p_{\text{start}}, n_{\text{start}})$~--- параметры
инициатора $\mathrm{ATG}$.
Все параметры неотрицательны и целочисленны.
Минимизируется суммарный вклад служебных кодонов
отдельно для нейтронов и протонов, как определено
в разделе~\ref{sec:parametrization}.
Сначала определим нейтронные параметры, затем протонные.

\subsubsection{Определение нейтронных параметров}
\label{sec:results-key1-neutrons}

\paragraph{Распределение служебных кодонов.}
Все четыре служебных кодона принадлежат Октету~II,
поэтому весь добавленный вклад приходится на этот октет.
Инициатор $\mathrm{ATG}$ и стоп-кодон $\mathrm{TAG}$
имеют $\mathrm{G}$ в третьей позиции и принадлежат
ветви Кето; стоп-кодоны $\mathrm{TAA}$ и $\mathrm{TGA}$
имеют $\mathrm{A}$ в третьей позиции и принадлежат
ветви Амино. Таким образом, ветвь Кето содержит один
стоп-кодон и инициатор, а ветвь Амино~--- два стоп-кодона.
То же распределение служебных позиций имеет место
в ветвях Сильные и Слабые. По тождествам групп-близнецов
(раздел~\ref{sec:twin-identities}) баланс Кето/Амино
одновременно обеспечивает баланс Сильные/Слабые.

\paragraph{Уравнение баланса для Кето и Амино.}
Полные нейтронные суммы складываются из сумм пула
без служебных кодонов
(раздел~\ref{sec:results-key1-trpile})
и искомых вкладов:
\begin{align*}
N_K(\mathrm{Keto}) &= 1813 + n_{\text{start}} + n_{\text{stop}}, \\
N_K(\mathrm{Amino}) &= 1776 + 2\,n_{\text{stop}}.
\end{align*}
Условие $N_K(\mathrm{Keto})=N_K(\mathrm{Amino})$ даёт
\[
n_{\text{stop}}-n_{\text{start}}=37.
\]
Правая часть равна разности нейтронных чисел триптофана и
изолейцина, к которой сведён дисбаланс пула. Баланс фиксирует
разность параметров, оставляя однопараметрическое семейство.
В основном выводе следующим условием служит делимость полной суммы.
Совместимость этого требования с выбором химической оси
разобрана в разделе~\ref{sec:results-key1-uniqueness}.

\paragraph{Делимость полной нейтронной суммы на~$37$.}
Для пула без служебных кодонов установлено
$N(\mathcal{C}_{\mathrm{SE}})=3589=97\cdot37$
(раздел~\ref{sec:results-sensepool}). При Key~0 вклад
инициатора равен нейтронному числу метионина, и
\[
N_0(\mathrm{All})=3589+69=3658=99\cdot37-5.
\]
Чтобы сохранить делимость на полной таблице, потребуем
\[
N_K(\mathrm{All})=
3589+n_{\text{start}}+3\,n_{\text{stop}}
\equiv0\pmod{37},
\]
то есть
\[
n_{\text{start}}+3\,n_{\text{stop}}\equiv0\pmod{37}.
\]

\paragraph{Решение.}
Подстановка $n_{\text{stop}} = n_{\text{start}} + 37$ из первого условия во второе дает:
\[
4\, n_{\text{start}} + 111 \equiv 0 \pmod{37}.
\]
Поскольку $111 = 3 \cdot 37 \equiv 0 \pmod{37}$ и $\gcd(4, 37) = 1$, это сводится к:
\[
n_{\text{start}} \equiv 0 \pmod{37},
\]
что дает семейство неотрицательных целочисленных решений
$n_{\text{start}} = 0, 37, 74, \ldots$.
Минимальное решение:
\[
n_{\text{start}} = 0,
\quad
n_{\text{stop}} = 37.
\]
Требование делимости не изменяет найденного минимума.
Баланс Кето/Амино при неотрицательных целочисленных
параметрах допускает $n_{\text{start}} = 0, 1, 2, \ldots$
и $n_{\text{stop}} = n_{\text{start}} + 37$. Делимость
сужает это семейство до
$n_{\text{start}} = 0, 37, 74, \ldots$, но минимальный
суммарный добавленный вклад в обоих случаях достигается
при $n_{\text{start}} = 0$ и $n_{\text{stop}} = 37$.
Для протонной компоненты требование делимости меняет
минимум. Структура семейств решений и роль каждого
из условий рассматриваются в
разделе~\ref{sec:results-key1-uniqueness}.

\paragraph{Исходные данные и результат вывода.}
В систему входят локальная разность $37$ и полная нейтронная
сумма пула $3589=97\cdot37$. Обе величины установлены до
назначения служебным кодонам. Коэффициент $4$ в объединённом
сравнении возникает при подстановке баланса в сумму вкладов
трёх стоп-кодонов и инициатора. Взаимная простота $4$ и $37$
позволяет сократить сравнение; это свойство уравнения, а
найденные параметры~--- его минимальное решение.

\paragraph{Совпадение с суммой Октета~I.}
После определения параметров полная нейтронная сумма равна
\[
N_1(\mathrm{All})=3589+0+3\cdot37
=3700=T(\mathrm{Octet~I}).
\]
Равенство с нуклонной суммой Октета~I не использовалось
как условие вывода. Оно выполняется уже при минимальном
нейтронном назначении, обеспечивающем только баланс.
На уровне Октета~II та же связь имеет вид
\[
N_1(\mathrm{Octet~II})=1889+0+3\cdot37
=2000=P(\mathrm{Octet~I}).
\]
Это переформулировка предыдущего равенства после вычитания
$N(\mathrm{Octet~I})=1700$, а не отдельное совпадение.
Баланс и тождества групп-близнецов дают
\[
\begin{aligned}
N_1(\mathrm{Keto})=N_1(\mathrm{Amino})
&=N_1(\mathrm{Strong})=N_1(\mathrm{Weak})\\
&=1850=\tfrac12\,T(\mathrm{Octet~I}).
\end{aligned}
\]
Здесь $1850$ получается делением найденной полной суммы пополам.
Дальнейшие следствия рассматриваются в
разделе~\ref{sec:results-key1-consequences}.

\subsubsection{Определение протонных параметров}
\label{sec:results-key1-protons}

Протонные параметры определяются отдельно от нейтронных
по тем же условиям: баланс Кето/Амино, делимость полной
суммы на~$37$ и минимальность. Равенство протонного
и нейтронного параметров стоп-кодона не предполагается.

\paragraph{Уравнение баланса для Кето и Амино.}
Из сумм пула без служебных кодонов
(раздел~\ref{sec:results-key1-trpile}) получаем
\begin{align*}
P_K(\mathrm{Keto}) &= 2108+p_{\text{start}}+p_{\text{stop}},\\
P_K(\mathrm{Amino}) &= 2072+2\,p_{\text{stop}}.
\end{align*}
Условие $P_K(\mathrm{Keto})=P_K(\mathrm{Amino})$ даёт
\[
p_{\text{stop}}-p_{\text{start}}=36.
\]
Правая часть равна $P(\mathrm{Trp})-P(\mathrm{Ile})$.
Чётность исключает значение $37$, но не определяет
фактическую разность $36$: она следует из протонных чисел
этих молекул (раздел~\ref{sec:results-proton-parity}).

\paragraph{Делимость полной протонной суммы на~$37$.}
Протонная сумма пула без служебных кодонов равна
\[
\begin{aligned}
P(\mathcal{C}_{\mathrm{SE}})
&=P(\mathrm{Octet~I})+P_{\mathrm{SE}}(\mathrm{Octet~II})\\
&=2000+2180=4180=113\cdot37-1.
\end{aligned}
\]
Это единичное отклонение от ближайшего кратного~$37$.
Для полной таблицы потребуем
\[
P_K(\mathrm{All})=
4180+p_{\text{start}}+3\,p_{\text{stop}}
\equiv0\pmod{37},
\]
то есть
\[
p_{\text{start}}+3\,p_{\text{stop}}\equiv1\pmod{37}.
\]

\paragraph{Решение.}
Подставляя $p_{\text{stop}}=p_{\text{start}}+36$, получаем
\[
4\,p_{\text{start}}+108\equiv1\pmod{37}.
\]
Поскольку $108=3\cdot37-3$ и $\gcd(4,37)=1$,
\[
4\,p_{\text{start}}\equiv4\pmod{37},
\qquad p_{\text{start}}\equiv1\pmod{37}.
\]
Неотрицательные целочисленные решения имеют
$p_{\text{start}}=1,38,75,\ldots$.
Минимальное решение:
\[
p_{\text{start}}=1,\qquad p_{\text{stop}}=37.
\]
Структура семейства и минимальность этого выбора
рассматриваются в разделе~\ref{sec:results-key1-uniqueness}.

\paragraph{Возникшее равенство параметров стопа.}
Протонная система использует локальную разность $36$ и
единичное отклонение полной суммы пула от кратного~$37$.
Как и в нейтронном выводе, коэффициент $4$ определяется
распределением служебных позиций и уравнением баланса.
Сокращение сравнения и выбор минимума дают параметры,
которые можно сопоставить с нейтронным решением
(раздел~\ref{sec:results-key1-neutrons}):
\[
p_{\text{stop}}=n_{\text{stop}}=37.
\]
Равенство протонного и нейтронного параметров не входило
в условия двух систем. Оно возникает при выборе их
минимальных неотрицательных целочисленных решений.

\paragraph{Макроскопические следствия.}
При Key~1 полная протонная сумма равна $4292$, а баланс
и тождества групп-близнецов дают
\[
\begin{aligned}
P_1(\mathrm{Keto})=P_1(\mathrm{Amino})
&=P_1(\mathrm{Strong})=P_1(\mathrm{Weak})\\
&=2146=\tfrac12\,T_1(\mathrm{Octet~II}).
\end{aligned}
\]
Соответствующее равенство полной таблицы имеет вид
\[
P_1(\mathrm{All})=T_1(\mathrm{Octet~II})=4292.
\]
Нуклонная сумма Октета~II здесь вычислена при Key~1 и
зависит от назначения служебным кодонам. Это равенство
следует из $N_1(\mathrm{Octet~II})=P(\mathrm{Octet~I})$
и определения $T=P+N$: оно выражает ту же связь между
полной таблицей и октетами, что и нейтронное совпадение.
Полный разбор приведён в
разделе~\ref{sec:results-key1-consequences}.

\subsubsection{Иерархия ограничений и выбор Key~1}
\label{sec:results-key1-uniqueness}

Системы разделов~\ref{sec:results-key1-neutrons}
и~\ref{sec:results-key1-protons} имеют общее построение.
При неотрицательных целочисленных параметрах выбор
назначения проходит три этапа:
\begin{enumerate}
\item Баланс Кето/Амино фиксирует разности
$n_{\text{stop}}-n_{\text{start}}=37$ и
$p_{\text{stop}}-p_{\text{start}}=36$.
Остаются два свободных параметра~--- значения инициатора
для нейтронов и протонов. Каждый элемент этого семейства
уравнивает четыре $32$-кодоновые группы Кето, Амино,
Сильные и Слабые по $P$, $N$, а следовательно, по $T$ и $\Delta$.
\item Делимость полных сумм на~$37$ ограничивает значения
инициатора до $n_{\text{start}}=37k$ и
$p_{\text{start}}=1+37m$, где $k,m$~--- неотрицательные
целые числа. Семейство по-прежнему имеет два свободных
целочисленных параметра, но шаг каждого увеличивается до~$37$.
\item Суммарные добавленные вклады равны $111+148k$
для нейтронов и $112+148m$ для протонов. Они строго
возрастают соответственно с $k$ и $m$, поэтому
единственный минимум достигается при $k=m=0$ и даёт Key~1:
$(p_{\text{stop}},n_{\text{stop}})=(37,37)$,
$(p_{\text{start}},n_{\text{start}})=(1,0)$.
\end{enumerate}
Единственность относится к минимуму при перечисленных
условиях: баланс и делимость допускают и большие назначения.
В частности, равенство протонного и нейтронного параметров
стоп-кодона выполняется при $k=m$, но их общее значение
равно $37$ только при $k=m=0$.

\paragraph{Вклад условия делимости.}
Сравним назначения с минимальным суммарным добавленным
вкладом, полученные отдельно для протонов и нейтронов при
одном балансе Кето/Амино и при дополнительном требовании
делимости на~$37$; в обоих случаях параметры
неотрицательны и целочисленны:
\[
\begin{aligned}
\text{только баланс:}\quad
&(p_{\text{stop}}, n_{\text{stop}}) = (36, 37), &
&(p_{\text{start}}, n_{\text{start}}) = (0, 0), \\
\text{баланс и делимость:}\quad
&(p_{\text{stop}}, n_{\text{stop}}) = (37, 37), &
&(p_{\text{start}}, n_{\text{start}}) = (1, 0).
\end{aligned}
\]
Требование делимости не изменяет минимальное нейтронное
назначение: нейтронная сумма полной таблицы равна $3700$
и кратна~$37$ уже при одном балансе. Минимальная протонная
сумма при одном балансе равна $4288 = 116 \cdot 37 - 4$;
прибавление по одному протону каждой из четырёх служебных
позиций увеличивает её до $4292 = 116 \cdot 37$. Баланс
при этом сохраняется, поскольку каждая половина оси
Кето/Амино содержит по две служебные позиции. Таким
образом, дополнительное требование делимости увеличивает
минимальный вклад на четыре протона и не требует
дополнительных нейтронов.

Та же прибавка устраняет разрыв по $\Delta$ на оси
Пурины/Пиримидины. У Пиримидинов $\Delta = 296$ не зависит
от параметризации (раздел~\ref{sec:results-pyrimidine}),
тогда как при минимальном назначении, обеспечивающем
только баланс Кето/Амино, у Пуринов $\Delta = 292$:
каждый стоп-кодон несёт $\Delta = 36 - 37 = -1$, что
смещает пуловое значение $295$
(раздел~\ref{sec:results-sensepool}) на три единицы вниз.
Все четыре служебных кодона принадлежат Пуринам, поэтому
прибавление по одному протону каждому из них уравнивает
значения обеих ветвей. Ниже рассматривается альтернативный
вывод, в котором это равенство задаётся вместе с балансом
Кето/Амино, а требование делимости не вводится.

\paragraph{Равенство параметров стопа как условие двух осей.}
Вместо делимости можно потребовать равенства $\Delta$
у Пуринов и Пиримидинов. Обозначим разности параметров
служебных кодонов через
$\Delta_{\text{stop}}=p_{\text{stop}}-n_{\text{stop}}$
и $\Delta_{\text{start}}=p_{\text{start}}-n_{\text{start}}$.
Баланс Кето/Амино даёт
\[
\Delta_{\text{start}}=\Delta_{\text{stop}}+1.
\]
Все служебные кодоны принадлежат Пуринам, а разность
между значениями $\Delta$ двух ветвей в пуле без служебных
кодонов равна
$\Delta_{\mathrm{SE}}(\mathrm{Pyrimidine})-
\Delta_{\mathrm{SE}}(\mathrm{Purine})=1$.
Поэтому дополнительный баланс требует
\[
\Delta_{\text{start}}+3\,\Delta_{\text{stop}}=1.
\]
Совместно эти уравнения дают
$\Delta_{\text{stop}}=0$ и $\Delta_{\text{start}}=1$.
Тем самым равенство $p_{\text{stop}}=n_{\text{stop}}$
и разность $p_{\text{start}}-n_{\text{start}}=1$
получаются без требования делимости.

Вместе с балансом Кето/Амино остаётся однопараметрическое семейство
$(p_{\text{stop}},n_{\text{stop}})=(37+t,37+t)$,
$(p_{\text{start}},n_{\text{start}})=(1+t,t)$ при целом $t\ge0$.
Вдоль него полная нейтронная сумма равна $3700+4t$,
а протонная~--- $4292+4t$. Минимальный суммарный вклад
достигается при $t=0$ и снова даёт Key~1.
Делимость обеих сумм на~$37$ выполняется только при
$t\equiv0\pmod{37}$; у минимального назначения она
является следствием, а не условием этого вывода.
Равенство стоповых параметров здесь обусловлено
согласованием разностей по двум осям: разности
$\delta P(\mathrm{Trp},\mathrm{Ile})-
\delta N(\mathrm{Trp},\mathrm{Ile})=-1$
(раздел~\ref{sec:results-key1-trpile}) и указанной выше
разности $\Delta$ между Пиримидинами и Пуринами.

\paragraph{Почему баланс по $P$ и $N$ выбран на оси Кето/Амино.}
Равенство протонных и нейтронных сумм Пуринов и Пиримидинов
потребовало бы других вкладов служебных кодонов:
\[
p_{\text{start}}+3\,p_{\text{stop}}=212,
\qquad n_{\text{start}}+3\,n_{\text{stop}}=211.
\]
Однако $212\equiv-10\pmod{37}$ и
$211\equiv-11\pmod{37}$, тогда как делимость полных сумм
требует соответственно остатков $1$ и $0$.
Поэтому такой баланс несовместим с делимостью на~$37$.
Среди трёх осей третьей позиции с этим требованием совместим
баланс по $P$ и $N$ на оси Кето/Амино и его структурная
копия на оси Сильные/Слабые. Баланс по $\Delta$ на оси
Пурины/Пиримидины не накладывает этих двух отдельных
равенств и совместим с найденным назначением.

\paragraph{Общее основание выводов.}
Оба вывода используют одинаковые назначения трём стоп-кодонам,
те же исходные суммы и баланс Кето/Амино. Дополнительное
условие различается: делимость полных сумм на~$37$ либо
равенство $\Delta$ у Пуринов и Пиримидинов. При
неотрицательности, целочисленности и минимальности оба
варианта дают Key~1. Это согласованность двух систем
условий с общим основанием, а не независимые подтверждения.

Полученное значение $1850=\tfrac12\,T(\mathrm{Octet~I})$
допускает обратную проверку нейтронных назначений:
$1850-1776=74=2\cdot37$ приходится на два стоп-кодона
ветви Амино, а $1850-1813=37$~--- на стоп-кодон и
инициатор ветви Кето. Отсюда снова получаются
$n_{\text{stop}}=37$ и $n_{\text{start}}=0$.
В этой проверке равенство с половиной суммы Октета~I
используется уже после его получения; оно не является
третьим независимым основанием вывода.

\subsection{Следствия при Key~1}
\label{sec:results-key1-consequences}

При Key~1 баланс Кето/Амино распространяется на четыре
$32$-кодоновые группы, а равенство нейтронной суммы полной
таблицы и нуклонной суммы Октета~I связывает несколько
уровней разбиения. В этом разделе рассматриваются эти
следствия, свойства делимости полученных сумм и остаточная
асимметрия Пуринов/Пиримидинов. Обзор групп при Key~1
приведён на Рис.~\ref{fig:groups-key1}.

\begin{figure}[!htbp]
  \centering
  \includegraphics[width=0.78\textwidth]{groups_key1.png}
\caption{Нуклонные суммы кодонных групп при Key~1:
каждому стоп-кодону назначено $(P,N)=(37,37)$,
инициатору $\mathrm{ATG}$~--- $(P,N)=(1,0)$.
Расположение групп и цветовой код делимости те же,
что на Рис.~\ref{fig:pool60}; сравнение с Key~0 дано
на Рис.~\ref{fig:groups-key0}.
Сгенерировано сопровождающим \texttt{visualization.html}.}
  \label{fig:groups-key1}
\end{figure}

\subsubsection{Выравнивание $32$-кодоновых групп химических осей}
\label{sec:results-key1-axisalignment}

При Key~1 группы Кето, Амино, Сильные и Слабые имеют
одинаковые значения четырёх нуклонных величин.
Для каждой из этих групп $G$:
\[
\begin{aligned}
T_1(G)&=3996, & P_1(G)&=2146,\\
N_1(G)&=1850, & \Delta_1(G)&=296.
\end{aligned}
\]
Баланс Кето/Амино вместе с тождествами групп-близнецов
обеспечивает равенство сумм всех четырёх групп
по $P$ и $N$, а определения $T=P+N$
и $\Delta=P-N$~--- по $T$ и $\Delta$.
Таким образом, четыре выровненные группы представляют
следствие выбранного назначения, а не четыре отдельных
подтверждения. Их суммы составляют половину сумм полной
таблицы. Все групповые значения приведены в
\texttt{key1\_nucleon-data.csv}.

Пурины и Пиримидины имеют то же значение $\Delta$, но
различаются по $P$, $N$ и $T$. Это различие разобрано
в разделе~\ref{sec:results-key1-purpyr}.

\subsubsection{Макроскопические тождества и распространение на химические половины осей}
\label{sec:results-key1-macroidentities}

Связь полной таблицы с октетами имеет три эквивалентные записи:
\begin{align*}
N_1(\mathrm{All})
&=T(\mathrm{Octet~I})=3700=100\cdot37,\\
P_1(\mathrm{All})
&=T_1(\mathrm{Octet~II})=4292=116\cdot37,\\
N_1(\mathrm{Octet~II})
&=P(\mathrm{Octet~I})=2000=2\cdot1000.
\end{align*}
Это одна связь, полученная в
разделе~\ref{sec:results-key1-neutrons} после определения
параметров. Вычитание $N(\mathrm{Octet~I})=1700$ из
первого равенства даёт третье. Прибавление к третьему
$P_1(\mathrm{Octet~II})$ даёт второе, поскольку
$T=P+N$. Обратные преобразования также выполняются.
Суммы Октета~I не зависят от параметризации; суммы
Октета~II включают назначения служебным кодонам.

Для каждой из четырёх выровненных групп $G$~--- Кето,
Амино, Сильные, Слабые~--- эта связь воспроизводится
на уровне её октетных пересечений:
\[
\begin{aligned}
N_1(G)&=T(G\cap\mathrm{Octet~I})=1850,\\
P_1(G)&=T_1(G\cap\mathrm{Octet~II})=2146,\\
N_1(G\cap\mathrm{Octet~II})
&=P(G\cap\mathrm{Octet~I})=1000.
\end{aligned}
\]
Каждое пересечение содержит $16$ кодонов. В Октете~I
каждая такая группа включает по два кодона каждого
четырёхкратно вырожденного семейства и потому несёт
половину сумм октета. В Октете~II равенство половин
по двум осям следует из баланса полной таблицы после
вычитания равных вкладов Октета~I. Поэтому приведённые
равенства представляют половинные записи той же связи.
Для Пуринов и Пиримидинов такое выравнивание Октета~II
не выполняется. Полный перечень равенств приведён в
\texttt{key1\_equalities.csv}.

\paragraph{Рациональный профиль.}
Полная таблица и четыре выровненные группы имеют
одинаковые отношения нуклонных величин при Key~1.
Обозначая любую из этих групп через $G$, получаем:
\[
\frac{P_1(G)}{T_1(G)}=\frac{29}{54},\qquad
\frac{N_1(G)}{T_1(G)}=\frac{25}{54},\qquad
\frac{\Delta_1(G)}{T_1(G)}=\frac{2}{27}.
\]
Соответственно, $P_1(G):N_1(G)=29:25$.
Это нормированная запись тех же сумм. Для каждой из
четырёх групп $2146=29\cdot2\cdot37$ одновременно
является протонной суммой группы и нуклонной суммой
её пересечения с Октетом~II; $1850=25\cdot2\cdot37$
связывает нейтронную сумму с пересечением с Октетом~I.
Профили отношений всех групп приведены в
\texttt{key1\_ratios.csv}.

\subsubsection{Делимость на \texorpdfstring{$37$}{37} и на \texorpdfstring{$999$}{999}}
\label{sec:results-key1-divisibility37and999}

\paragraph{Делимость на $37$.}
В основном выводе Key~1 делимость полных сумм по $P$ и
$N$ задана явно. Для нейтронов она выполняется уже
при минимальном назначении, обеспечивающем только
баланс; для протонов меняет минимум
(раздел~\ref{sec:results-key1-uniqueness}).
Полученные суммы равны
\[
N_1(\mathrm{All})=100\cdot37,\qquad
P_1(\mathrm{All})=116\cdot37.
\]
Баланс делит каждую сумму на две равные целочисленные
части. Поскольку $37$ нечётно, обе части также кратны
$37$. Для каждой из четырёх выровненных групп $G$ имеем
\[
\begin{aligned}
T_1(G)&=108\cdot37, & P_1(G)&=58\cdot37,\\
N_1(G)&=50\cdot37, & \Delta_1(G)&=8\cdot37.
\end{aligned}
\]
Делимость $T$ и $\Delta$ следует сложением и вычитанием
$P$ и $N$; в частности,
$\Delta_1(\mathrm{All})=4292-3700=592=16\cdot37$.

Для нуклонной суммы действует и более общее свойство:
\[
T_1(G)\equiv T_0(G)\pmod{37}
\]
для любой группы кодонов $G$. При переходе от Key~0
к Key~1 вклад каждого стоп-кодона в $T$ возрастает
на $74=2\cdot37$, а вклад $\mathrm{ATG}$ меняется
на $1-149=-148=-4\cdot37$. Поэтому делимость $T$
на $37$ сохраняется при этом переходе независимо
от состава группы. Перечень делимостей при Key~1
приведён в \texttt{key1\_divisibility-37.csv}.

\paragraph{Делимость на $999$.}
Полная нуклонная сумма имеет вид
\[
T_1(\mathrm{All})=7992=8\cdot999=6^3\cdot37.
\]
В отличие от делимости на $37$, делимость на
$999=27\cdot37$ не задавалась отдельным условием.
Она характеризует полученное минимальное решение.
Нуклонная сумма каждой из четырёх выровненных половин
при Key~1 равна $3996=4\cdot999$; это следствие деления
полной суммы пополам.

\subsubsection{\texorpdfstring{$\Delta$}{Delta}-инвариант во всех шести $32$-кодоновых группах}
\label{sec:results-key1-deltainvariant}

При Key~1 все шесть $32$-кодоновых групп третьей позиции
имеют одинаковую протонно-нейтронную разность:
\[
\begin{aligned}
\Delta_1(\mathrm{Keto})=\Delta_1(\mathrm{Amino})
&=\Delta_1(\mathrm{Strong})=\Delta_1(\mathrm{Weak})\\
&=\Delta_1(\mathrm{Purine})=\Delta(\mathrm{Pyrimidine})=296.
\end{aligned}
\]
Для первых четырёх групп равенство следует из баланса
по $P$ и $N$. Равенство Пуринов и Пиримидинов по
$\Delta$ в основном выводе не задаётся отдельно;
в альтернативном двухосевом выводе оно служит
условием вместо делимости
(раздел~\ref{sec:results-key1-uniqueness}).
Поэтому его появление в этих двух выводах имеет
разный логический статус.

Общая разность воспроизводится на нижних уровнях
разбиения. Каждая восьмикодоновая группа Октета~I,
выделенная одним нуклеотидом третьей позиции, имеет
$\Delta=75$ независимо от параметризации
(раздел~\ref{sec:results-octetI}). При Key~1 каждая
соответствующая группа Октета~II имеет $\Delta=73$:
\[
\Delta(G_X\cap\mathrm{Octet~I})=75,\qquad
\Delta_1(G_X\cap\mathrm{Octet~II})=73
\qquad (|X|=1).
\]
В пиримидиновых группах значение $73$ также не зависит
от параметризации (раздел~\ref{sec:results-pyrimidine});
в пуриновых оно достигается при найденном назначении.
Сложение двух восьмикодоновых групп внутри каждого
октета даёт
\[
\Delta(G_X\cap\mathrm{Octet~I})=150,\qquad
\Delta_1(G_X\cap\mathrm{Octet~II})=146
\qquad (|X|=2).
\]
Сложение по октетам даёт для каждой однонуклеотидной
группы полной таблицы
\[
\Delta_1(G_X)=75+73=148=4\cdot37
\qquad (|X|=1).
\]
При Key~0 группа $\{\mathrm{G}\}$ имела
$\Delta_0(\{\mathrm{G}\})=158$, тогда как остальные
три уже имели $148$
(раздел~\ref{sec:results-key0-delta}). При Key~1
значения всех четырёх совпадают. Оба способа сложения
завершаются на уровне $32$ кодонов:
\[
296=2\cdot148=150+146.
\]
Таким образом, каскад описывает одну согласованность
на нескольких уровнях. Она сохраняется и у неминимальных
назначений двухосевого семейства из
раздела~\ref{sec:results-key1-uniqueness}: увеличение
обоих параметров на одну и ту же величину не меняет
разность $\Delta$ ни у одного служебного кодона.

\subsubsection{Асимметрия Пуринов/Пиримидинов}
\label{sec:results-key1-purpyr}

При Key~1 ось Пурины/Пиримидины сохраняет следующие
разности между суммами двух половин:
\begin{align*}
N(\mathrm{Pyrimidine})-N_1(\mathrm{Purine})&=100,\\
P(\mathrm{Pyrimidine})-P_1(\mathrm{Purine})&=100,\\
T(\mathrm{Pyrimidine})-T_1(\mathrm{Purine})&=200,\\
\Delta(\mathrm{Pyrimidine})-\Delta_1(\mathrm{Purine})&=0.
\end{align*}
Вся асимметрия приходится на Октет~II: в Октете~I
пуриновые и пиримидиновые половины имеют равные
суммы вследствие четырёхкратной вырожденности семейств.

В пуле без служебных кодонов разности Пиримидины--Пурины
составляют $212$ по $P$ и $211$ по $N$
(раздел~\ref{sec:results-sensepool}). Назначение Key~1
добавляет в пуриновую половину $112$ протонов и
$111$ нейтронов, поэтому
\[
212-112=211-111=100.
\]
Равенство двух сдвигов эквивалентно равенству $\Delta$
у Пуринов и Пиримидинов. Их общая величина $100$
получается из исходных сумм после назначения параметров;
она не задавалась условием вывода.

\paragraph{Равенства на однонуклеотидном уровне.}
Равенство сумм групп $\{\mathrm{C}\}$ и
$\{\mathrm{T}\}$ не зависит от параметризации и
следует из пиримидиновой синонимии
(раздел~\ref{sec:twin-identities}). Поскольку Кето
объединяет $\{\mathrm{G}\}$ и $\{\mathrm{T}\}$,
а Амино~--- $\{\mathrm{A}\}$ и $\{\mathrm{C}\}$,
баланс по $P$ и $N$ даёт
\[
P_1(\{\mathrm{G}\})=P_1(\{\mathrm{A}\}),\qquad
N_1(\{\mathrm{G}\})=N_1(\{\mathrm{A}\}).
\]
Это переформулировка баланса, верная при любом
назначении, которое его обеспечивает. При Key~1
две пары групп имеют значения
\[
\begin{aligned}
\bigl(T(\{\mathrm{C}\}),P(\{\mathrm{C}\}),N(\{\mathrm{C}\})\bigr)
&=(2048,1098,950),\\
\bigl(T_1(\{\mathrm{G}\}),P_1(\{\mathrm{G}\}),N_1(\{\mathrm{G}\})\bigr)
&=(1948,1048,900),
\end{aligned}
\]
с теми же значениями у $\{\mathrm{T}\}$ и
$\{\mathrm{A}\}$ соответственно. Разность между
парами равна $50$ по $P$ и $N$ и $100$ по $T$;
удвоение воспроизводит разности двух $32$-кодоновых
половин. После вычитания одинаковых вкладов Октета~I
в каждой из групп
$\{\mathrm{G}\}\cap\mathrm{Octet~II}$ и
$\{\mathrm{A}\}\cap\mathrm{Octet~II}$
при Key~1 нуклонная сумма равна $1023$,
протонная~--- $548$, нейтронная~--- $475$.

\paragraph{Внутри Октета~II на уровне $16$ кодонов.}
Пуриновая и пиримидиновая половины Октета~II имеют
нейтронные суммы $950$ и $1050$ и протонные суммы
$1096$ и $1196$ соответственно. Они отклоняются на $\pm50$ от средних значений:
$1000$ для нейтронной суммы и $1146$ для протонной.
Такие суммы при Key~1 имеют половины этого октета
вдоль осей Кето/Амино и Сильные/Слабые. Это ещё одна запись
той же асимметрии: две половины симметричны относительно
среднего, но не равны друг другу.

\paragraph{Предел выравнивания.}
Остаточная асимметрия не означает невозможности полного
баланса трёх осей. В двухосевом семействе из
раздела~\ref{sec:results-key1-uniqueness}, где стоп-кодоны
имеют параметры $(37+t,37+t)$, а инициатор~--- $(1+t,t)$,
оба сдвига Пиримидины--Пурины равны $100-4t$.
При $t=25$ назначение стоп-кодонам $(62,62)$ и
инициатору $(26,25)$ уравнивает все шесть половин
по $P$ и $N$. Однако полные суммы становятся равными
$P=4392$ и $N=3800$ и не делятся на $37$;
нейтронное равенство с $T(\mathrm{Octet~I})=3700$
также теряется. Key~1 фиксирует минимальное согласование
баланса с делимостью, сохраняя указанную асимметрию.

\subsection{Синтез}
\label{sec:results-synthesis}

Рассмотренные закономерности связывают два уровня организации
стандартного генетического кода: структуру кодонных групп и
нуклонный состав кодируемых аминокислот. Разбиения опираются на
химические свойства оснований, пиримидиновую синонимию и
октеты Румера. В принятом представлении свободных
нейтральных аминокислот числа протонов и нейтронов однозначно
определяются молекулярным составом, а групповые суммы учитывают
число кодонов каждой аминокислоты. Выбор третьей позиции,
молекулярного представления и модуля~$37$ задаёт схему анализа;
наблюдаемые числовые значения из этих выборов сами по себе
не следуют.

Уже без назначения значений служебным кодонам структурные
разбиения обнаруживают арифметическое согласование.
В Октете~I протонная и нейтронная суммы равны $2000$ и $1700$,
а их сумма $3700=100\cdot37$ соединяет десятичный масштаб
с модулем~$37$. Структура октета обеспечивает точное
масштабирование этих величин в его группах, но не определяет
их исходные значения. На уровне всего пула без служебных
кодонов нейтронная сумма равна $3589=97\cdot37$, а протонная
составляет $4180=113\cdot37-1$
(раздел~\ref{sec:results-independent}).

Локальная асимметрия согласуется с этим глобальным рисунком.
Различие Кето и Амино сводится к паре
триптофан--изолейцин, дающей разности $37$ по нейтронам
и $36$ по протонам. Среди $190$ пар канонических аминокислот такое сочетание разностей встречается только у пар триптофан--изолейцин и
триптофан--лейцин. Совпадение объясняется тем, что
изолейцин и лейцин являются изомерами и обладают
одинаковыми нуклонными числами
(раздел~\ref{sec:results-key1-trpile}).
Чётность протонных чисел всех канонических аминокислот
исключает протонную разность $37$, но не определяет,
насколько фактическая разность должна от неё отклоняться.
Значение $36$ оказывается ближайшим допустимым снизу,
и разность между протонным и нейтронным дисбалансами
составляет ровно единицу
(раздел~\ref{sec:results-proton-parity}).

Единичные отклонения встречаются на разных уровнях структуры и имеют различное арифметическое происхождение.
В Октете~I значения $75=2\cdot37+1$ возникают при
делении его разности $\Delta=300$ между четырьмя равными
восьмикодоновыми группами. В пиримидиновых группах
Октета~II значение противоположно смещено:
$73=2\cdot37-1$. При объединении октетов отклонения
компенсируются: $75+73=148=4\cdot37$.
В пуле без служебных кодонов значение $\Delta=73$
имеет также ветвь $\{\mathrm{A}\}$ Октета~II,
что уже не следует из пиримидиновой синонимии.
Отличается только ветвь $\{\mathrm{G}\}$ со значением
$72$; в ней находится исключённая позиция $\mathrm{ATG}$
(раздел~\ref{sec:results-key0-delta}).
Ещё одно согласование проявляется при Key~0:
полное число нуклонов метионина $149=4\cdot37+1$
компенсирует остаток $-1$ нуклонной суммы пула.
Таким образом, единичные отклонения имеют несколько
источников, а их повторения распространяются через
структурные связи между группами.

Расположение служебных кодонов позволяет согласовать
локальный баланс с глобальными суммами.
В Кето находятся один стоп-кодон и $\mathrm{ATG}$,
в Амино~--- два стоп-кодона. При одинаковом назначении
стоп-кодонам минимальная неотрицательная целочисленная
компенсация нейтронного дисбаланса даёт каждому из них
$37$ нейтронов, а $\mathrm{ATG}$~--- ноль. После этого
\[
N_1(\mathrm{All})=3589+3\cdot37
               =3700=T(\mathrm{Octet~I}).
\]
Совпадение с суммой Октета~I не использовалось при выборе
компенсации. Здесь согласуются величина локального дисбаланса,
число и расположение стоп-кодонов и сумма
структурного Октета~I.

На протонной стороне переход от минимального баланса
к минимальному балансовому назначению с делимостью
на~$37$ требует четырёх дополнительных протонов~---
по одному каждой служебной позиции.
Их распределение поровну между Кето и Амино сохраняет
баланс. Принадлежность всех четырёх позиций Пуринам
позволяет той же прибавкой устранить разрыв по $\Delta$
между Пуринами и Пиримидинами, сохраняя исходные
пиримидиновые суммы. Получается Key~1:
$(P,N)=(37,37)$ для каждого стоп-кодона и $(1,0)$
для $\mathrm{ATG}$
(раздел~\ref{sec:results-key1-derivation}).
Его простота выражает совместимость исходных чисел
с расположением служебных позиций; единственность
относится к принятым условиям и критерию минимальности.

Полученное назначение объединяет эти отношения
в полной таблице. Несколько макроскопических равенств
выражают одну связь с Октетом~I; равенство $\Delta$
совместимо с одинаковыми сдвигами протонных и нейтронных
сумм Пуринов и Пиримидинов на $100$.
Нуклонная сумма $7992=8\cdot999$ дополняет числовой
профиль, хотя делимость на~$999$ не задавалась
отдельным условием
(раздел~\ref{sec:results-key1-consequences}).
Содержание результата состоит в согласовании молекулярных
чисел, структурных разбиений и формальных назначений;
повторные проявления одной связи не являются отдельными
свидетельствами. Специфичность этого согласования
оценивается далее сравнением с альтернативными кодами,
детерминированными контролями и нулевой моделью
(раздел~\ref{sec:discussion}).

\section{Обсуждение}
\label{sec:discussion}

В данном разделе исследуется, насколько описанные выше
закономерности специфичны для стандартного кода и его
структуры, а не являются артефактами проведённого анализа.
Раздел~\ref{sec:discussion-trpile-robustness} сравнивает
стандартный код с альтернативными кодами NCBI по
нейтронному дефициту и протонному остатку пула без
служебных кодонов, локализации асимметрии в паре
триптофан--изолейцин и результатам повторного назначения
служебным кодонам;
раздел~\ref{sec:discussion-controls} ставит вопрос о том,
выделяются ли третья кодонная позиция и модуль~$37$ среди
других позиций и делителей; а
раздел~\ref{sec:discussion-montecarlo} оценивает
концентрацию значений на решётке с шагом~$37$ относительно
нулевой модели случайных кодов. Раздел~\ref{sec:discussion-prior-work}
соотносит полученные результаты с предшествующими исследованиями, а
раздел~\ref{sec:discussion-limitations} излагает
ограничения подхода и перспективы дальнейшей работы.

\subsection{Специфичность среди альтернативных генетических кодов}
\label{sec:discussion-trpile-robustness}
\label{sec:results-key1-ncbi}

Сравнение $27$ таблиц трансляции из реестра NCBI
показывает, какие арифметические свойства стандартного
кода встречаются в других генетических кодах.
Для пула без служебных кодонов сравниваются
нейтронный дефицит и протонный остаток
(раздел~\ref{sec:results-sensepool}),
а также разности нейтронных и протонных сумм
между Кето и Амино. Для таблиц, воспроизводящих
стандартные значения этих разностей, проверяется
сохранение их локализации в паре
триптофан--изолейцин
(раздел~\ref{sec:results-key1-trpile}).
Отдельно к каждой таблице применяются одни и те же
условия баланса, делимости и минимальности
(раздел~\ref{sec:results-key1-derivation}), и
сравниваются полученные назначения служебным кодонам
и групповые суммы.

Для всех сравнений используются две модели.
В Модели~1 исключаются
$\mathrm{ATG}$ и все кодоны, перечисленные в реестре
как стоп-кодоны; в Модели~2 исключаются $\mathrm{ATG}$
и только постоянные стоп-кодоны. Контекстно-зависимые
стоп-кодоны в Модели~2 учитываются с указанными
аминокислотными назначениями. Различие затрагивает
таблицы $27$, $28$ и $31$. Исключение именно
$\mathrm{ATG}$ при сохранении элонгационных назначений
альтернативных старт-кодонов является общим соглашением
сравниваемых моделей
(раздел~\ref{sec:genetic-code}).

\paragraph{Нейтронный дефицит и протонный остаток.}
Для стандартного кода нейтронная сумма пула без
служебных кодонов меньше нуклонной суммы Октета~I на
$3700-3589=111=3\cdot37$, а протонная сумма имеет
остаток $-1$ по модулю~$37$. Первый тест проверяет
воспроизведение этих двух значений в других таблицах.

В \texttt{deficit\_models\_analysis.csv} нейтронный
дефицит вычисляется относительно трёх опорных сумм:
константы $3700$ стандартного кода; нуклонной суммы
структурного Октета~I, то есть тех же восьми кодонных
семейств, что и в стандартном коде, но при
аминокислотных назначениях проверяемой таблицы;
нуклонной суммы функционального Октета~I, то есть всех
полностью четырёхкратно вырожденных семейств самой
проверяемой таблицы (раздел~\ref{sec:datasets}).
Все три дефицита определяются для каждой из двух
моделей пула, что даёт шесть вариантов сравнения.
Протонный остаток вычисляется по модулю~$37$
непосредственно из протонной суммы пула без служебных
кодонов отдельно для каждой модели и не требует
опорной суммы.

Нейтронный дефицит $111$ при любой из трёх опорных
сумм и протонный остаток $-1$ по модулю~$37$
по отдельности выделяют одни и те же таблицы:
$1$, $11$ и $28$ в Модели~1 и только $1$ и $11$
в Модели~2. В остальных $24$ таблицах ни один
из двух признаков не воспроизводится ни при одной
модели.

Таблица~11 имеет те же аминокислотные назначения
и стоп-кодоны, что и стандартная. В таблице~28
исключение трёх контекстно-зависимых стоп-кодонов
в Модели~1 оставляет тот же пул; их возвращение
в Модели~2 нарушает оба совпадения. Таким образом,
совпавшие строки представляют одну и ту же
арифметику исходного пула, а не независимые
повторения результата.

\paragraph{Сохранение асимметрии Кето/Амино и её локализации.}
Второй тест проверяет, воспроизводятся ли разности
нейтронных и протонных сумм Кето минус Амино,
равные в стандартном коде $37$ и $36$ соответственно.
Значения для всех таблиц и обеих моделей приведены
в \texttt{keto\_amino\_balance\_models.csv}.
Эта пара разностей воспроизводится в таблицах
$1$, $6$, $11$, $29$ и $30$ при обеих моделях,
а в таблицах $27$ и $28$~--- только при Модели~1.
Само совпадение разностей ещё не устанавливает их
локализацию; её сохранение проверяется по назначениям
кодонов.

В таблицах $6$, $29$ и $30$ кодоны $\mathrm{TAA}$
и $\mathrm{TAG}$ переназначены одной и той же
аминокислоте: глутамину, тирозину или глутаминовой
кислоте соответственно. В семействе $\mathrm{TA}$
кодон $\mathrm{TAG}$ принадлежит ветви Кето,
а $\mathrm{TAA}$~--- ветви Амино, поэтому добавленные
вклады сокращаются в разности. Других изменений
в этих таблицах нет: $\mathrm{TGG}$ и все три кодона
изолейцина сохраняют стандартные назначения, и
локализация в паре триптофан--изолейцин сохраняется.
Тот же механизм действует в таблице~27 при Модели~1:
$\mathrm{TAA}$ и $\mathrm{TAG}$ кодируют глутамин,
а контекстно-зависимый $\mathrm{TGA}$ исключён.
В Модели~2 его вклад триптофана нарушает совпадение.
Таблицы $11$ и $28$ при Модели~1 сохраняют исходный
пул стандартного кода; таблица~11 сохраняет его и
при Модели~2.

В Модели~1 пару разностей воспроизводят стандартная
таблица, пять альтернативных ядерных таблиц
и таблица~11, используемая бактериями, археями
и пластидами. Ни одна из $13$ митохондриальных
таблиц её не воспроизводит. Нарушение связано
не только с переназначением $\mathrm{ATA}$ или
$\mathrm{TGA}$: в таблицах $16$, $22$ и $23$ оба
кодона сохраняют стандартные назначения, тогда как
в остальных десяти изменён хотя бы один из них.
В таблице~16 $\mathrm{TAG}$ кодирует лейцин;
в таблице~22 к этому добавляется превращение
$\mathrm{TCA}$ в стоп-кодон; в таблице~23 стоп-кодоном
становится $\mathrm{TTA}$. После этих переназначений в разности сумм Кето
и Амино появляются дополнительные вклады, поэтому
она уже не сводится к разности нуклонных чисел
триптофана и изолейцина.

\paragraph{Сохранение разностей и делимость сумм.}
В стандартном пуле без служебных кодонов нейтронные
суммы Кето и Амино кратны~$37$. Протонная сумма
Амино также кратна~$37$, а протонная сумма Кето
имеет остаток $-1$ по модулю~$37$.
Этот набор остатков воспроизводят ровно те же
таблицы, что и стандартные значения нейтронного
дефицита и протонного остатка в первом тесте:
$1$, $11$ и $28$ в Модели~1 и $1$, $11$
в Модели~2.

В таблицах $6$, $27$, $29$ и $30$ симметричное
переназначение $\mathrm{TAA}$ и $\mathrm{TAG}$
вносит одинаковый вклад в обе половины отдельно
по протонам и нейтронам; эти вклады не кратны~$37$.
В тех моделях, где сохраняются стандартные
разности, ни одна из четырёх сумм не кратна~$37$:
нейтронные остатки одинаковы и ненулевые,
а протонный остаток Кето на единицу меньше
протонного остатка Амино.

Таким образом, локализация асимметрии в паре
триптофан--изолейцин сохраняется и при симметричном
переназначении $\mathrm{TAA}$ и $\mathrm{TAG}$.
Однако одинаковые прибавки к обеим половинам,
сохраняя их разности, меняют остатки самих сумм
по модулю~$37$. Поэтому сохранение этой локальной
асимметрии ещё не означает воспроизведения
арифметических закономерностей исходного пула
стандартного кода.

\paragraph{Сравнение минимальных назначений.}
Для каждой таблицы и обеих моделей набор
\texttt{service\_codon\_balance\_comparison.csv}
сопоставляет два минимальных назначения: при одном
балансе Кето/Амино и при добавлении глобальной
делимости на $37$. В каждой процедуре всем
стоп-кодонам принятой модели назначается одинаковое
значение, а $\mathrm{ATG}$~--- отдельное;
минимизируется суммарный добавленный вклад при
неотрицательных целочисленных параметрах отдельно
по $P$ и $N$. Параметры выводятся заново для каждой
таблицы.

Баланс с делимостью имеет совместное решение
по протонам и нейтронам не для всех таблиц:
в Модели~1 его нет у таблиц $14$, $25$ и $33$,
а в Модели~2~--- также у $27$, $28$ и $31$.
При одном балансе совместного решения нет только
у таблицы~25. У таблиц $14$ и $33$ минимальный
суммарный вклад определён отдельно по протонам
и нейтронам, но его распределение между служебными
кодонами не определяется однозначно.

При одном балансе минимальная нейтронная компенсация
приводит к равенству нейтронной суммы полной таблицы
и нуклонной суммы её структурного Октета~I только
в таблицах $1$, $11$ и $28$ при Модели~1 и
$1$, $11$ при Модели~2. Ни это равенство, ни
делимость на $37$ в процедуре не задаются.
Для стандартного кода получаются $37$ нейтронов
на каждый стоп-кодон и нулевой вклад $\mathrm{ATG}$;
сумма равна $3589+3\cdot37=3700$.

При дополнительном требовании делимости результат
расширяется: равенство с нуклонной суммой структурного
Октета~I выполняется также у таблиц
$2$, $9$, $10$, $13$, $21$ и $24$ при обеих моделях.
Среди таблиц с совместным решением баланса
с делимостью это равенство достигается у $9$ из $24$
таблиц в Модели~1 и у $8$ из $21$ в Модели~2.
После минимального назначения по обеим компонентам
у этих шести таблиц совпадают со стандартным кодом
при Key~1 все четыре величины во всех $33$ группах.
Однако сами назначения отличаются. Например, для
таблицы~2 получаются $(P,N)=(26,24)$ на стоп-кодон
и $(80,69)$ на $\mathrm{ATG}$, а для таблицы~10~---
$(10,17)$ и $(28,20)$ соответственно. Таблица~10
отличается от стандартной только чтением
$\mathrm{TGA}$ как цистеина; у остальных пяти
изменены также другие аминокислотные назначения.

У таблиц $12$ и $26$ после этой процедуры протонная
и нейтронная суммы также совпадают со стандартными
значениями $4292$ и $3700$, но полный профиль групп
отличается. Следовательно, совпадение полных сумм,
включая нуклонную сумму $7992=8\cdot999$, ещё
не означает совпадения сумм отдельных групп.

То же минимальное назначение, что у стандартного
кода,~--- $(37,37)$ каждому стоп-кодону и $(1,0)$
инициатору $\mathrm{ATG}$~--- получается только
для таблиц $1$, $11$ и $28$ в Модели~1 и
$1$, $11$ в Модели~2. Это выделенность формы
назначения, а не величины самого добавленного
вклада: таблицы $9$, $10$, $13$, $21$ и $24$
воспроизводят полный профиль при меньшем суммарном
добавленном вкладе как по $P$, так и по $N$.

Дополнительный вклад, необходимый для выполнения
делимости после достижения минимального баланса,
у стандартного кода составляет четыре протона
и ноль нейтронов. Среди таблиц с совместным решением
это наименьшая прибавка отдельно по каждой компоненте.
Те же значения имеют таблица~11 в обеих моделях
и таблица~28 в Модели~1. У остальных таблиц
с совместным решением обе прибавки больше.

Таким образом, стандартные значения нейтронного
дефицита и протонного остатка воспроизводятся
только у таблиц с тем же пулом без служебных кодонов.
Локализация асимметрии в паре триптофан--изолейцин
сохраняется и при некоторых изменениях этого пула.
После назначения служебным кодонам совпадения
становятся шире: другие коды могут воспроизводить
весь профиль групп стандартного кода, но при иных
значениях параметров. Поэтому совпадение итоговых
сумм само по себе не означает совпадения исходной
арифметики или назначения служебным кодонам.

Эти выводы относятся к рассмотренным таблицам
реестра. Число совпадений не даёт оценки вероятности,
поскольку таблицы не являются независимой случайной
выборкой генетических кодов.

\subsection{Специфичность третьей позиции кодона, модуля~\texorpdfstring{$37$}{37} и молекулярного представления}
\label{sec:discussion-controls}

Третья позиция кодона, модуль~$37$ и представление свободных
нейтральных аминокислот задают схему анализа
(разделы~\ref{sec:partitioning}, \ref{sec:criteria}
и~\ref{sec:composition}). Ниже проверяется, какие свойства
сохраняются при изменении этих выборов. Позиционный контроль
сопоставляет делимость сумм по одним и тем же химическим осям,
скан простых~--- распространённость делимости по разным модулям,
а контроль представления~--- устойчивость абсолютных сумм
и разностей. Затем рассматривается дополнительное
апостериорное наблюдение об асимметрии Кето/Амино
по трём позициям.

\paragraph{Позиция кодона.}
Позиционный контроль рассмотрим для пула без служебных
кодонов: его суммы не зависят от назначений служебным кодонам.
Применим три химические оси поочерёдно к каждой позиции.
На третьей позиции
у групп Амино и Слабые все четыре величины кратны~$37$;
у Кето и Сильных кратны~$37$ нейтронные суммы,
а у Пиримидинов~--- разность~$\Delta$
(раздел~\ref{sec:results-sensepool}).
У Кето и Сильных протонная и нуклонная суммы при этом
находятся на единицу ниже ближайшего кратного~$37$:
\[
\begin{aligned}
P_{\mathrm{SE}}(\mathrm{Keto})
  = P_{\mathrm{SE}}(\mathrm{Strong}) &= 2108=57\cdot37-1,\\
T_{\mathrm{SE}}(\mathrm{Keto})
  = T_{\mathrm{SE}}(\mathrm{Strong}) &= 3921=106\cdot37-1.
\end{aligned}
\]
Поскольку нейтронные суммы кратны~$37$, тождества
$T=P+N$ и $\Delta=P-N$ переносят остаток протонной суммы
на $T$ и $\Delta$.
На первой позиции
делятся только нейтронные суммы Кето и Амино.
При известной делимости полной суммы пула
$N(\mathcal C_{\mathrm{SE}})=3589=97\cdot37$
делимость одной из этих половин влечёт делимость другой,
поэтому два совпадения не являются независимыми.
На второй позиции ни одна из четырёх величин
ни одной из шести групп не делится на~$37$.

Химические оси делят полную таблицу на половины по~$32$
кодона, но после исключения служебных позиций размеры
могут различаться. В частности, Кето и Амино первой позиции
содержат $29$ и $31$ кодон, а третьей~--- по~$30$.
Таким образом, в рассмотренном пуле делимости групповых
сумм сосредоточены на третьей позиции; перенос тех же
химических классификаций на другие позиции не воспроизводит
этот набор делимостей. Значения для всех трёх позиций
в пуле без служебных кодонов (обозначен в файле
значением \texttt{sense}), при Key~0 и при Key~1
приведены в \texttt{position\_axis\_analysis.csv}.

\paragraph{Выбор модуля.}
Скан охватывает все простые числа до~$97$ и три набора
групповых величин: значения $T$, $P$, $N$ и $\Delta$
тех $17$ групп из~$33$, которые не содержат служебных
кодонов; значения всех $33$ групп при Key~0; значения
тех же групп при Key~1. В файле наборы обозначены
значениями \texttt{param\_independent}, \texttt{key0}
и \texttt{key1}.
В первом наборе назначения служебным кодонам не участвуют
в суммах, во втором
делимость на~$37$ не задаётся, а в основном выводе
Key~1 она входит в условия для полных сумм $P$ и $N$.

Для каждого простого числа $p$ подсчитывается доля различных
ненулевых значений, которые оно делит. Повторяющиеся
численные значения учитываются однократно. Эту долю
сопоставляем с ориентиром $1/p$, соответствующим
равномерному распределению остатков. Во всех трёх наборах
наибольшее отношение к этому ориентиру имеет~$37$:
оно делит $5$ из $28$ различных значений в первом наборе,
$13$ из $70$ при Key~0 и $10$ из $47$ при Key~1.
Соответствующие отношения равны приблизительно
$6{,}61$, $6{,}87$ и $7{,}87$.
Для $2$ и $5$ они меньше трёх во всех трёх наборах;
делимость на эти простые связана соответственно
с чётностью протонных чисел и с десятичной круглостью,
рассмотренными в разделах~\ref{sec:results-proton-parity}
и~\ref{sec:results-synthesis}.

Отношение к $1/p$ является описательной характеристикой:
скан не задаёт нулевой модели генетических кодов
и не даёт $p$-значения. Если из $m$ различных значений
на $p$ делятся $k$, отношение равно $pk/m$:
оно не превосходит $p$ и изменяется с шагом $p/m$.
Таким образом, диапазон и шаг показателя зависят
от размера модуля. Удаление одинаковых значений
не устраняет зависимостей, обусловленных сложением
и масштабированием групповых сумм.
Область сравнения ограничена простыми числами до~$97$.
Степени простых задают дополнительные условия делимости
и не входят в это сопоставление. Отдельно установлено,
что ни одно значение в трёх наборах не делится на~$37^2$.

Дополнительно каждое положительное значение приводится
к нечётному ядру удалением всех множителей~$2$.
Так устраняются повторы вида $v$, $2v$ и $4v$,
возникающие, в частности, при точном масштабировании
групп Октета~I (раздел~\ref{sec:results-octetI}).
При Key~0 число~$37$ делит десять различных нечётных
ядер, тогда как кратные каждого из чисел $73$ и $61$
сводятся к одному ядру. Для $73$ это само значение
$73=2\cdot37-1$ из каскада~$\Delta$ Октета~II;
для $61$~--- значение $549=9\cdot61$, к которому
сводятся протонные суммы пиримидиновых групп:
$2196=4\cdot549$ и $1098=2\cdot549$.
Этот подсчёт не устанавливает независимости
оставшихся значений.
Полные счётчики приведены в
\texttt{prime\_divisibility\_scan.csv}.
Первое место~$37$ по выбранной метрике наблюдается уже
в наборе для групп без служебных кодонов и при Key~0;
оно не обусловлено назначениями Key~1.

\paragraph{Зависимость от молекулярного представления.}
Сопоставим принятое представление свободной нейтральной
аминокислоты с пептидным остатком и цвиттер-ионом,
сохраняя изотопное соглашение
раздела~\ref{sec:composition}. В модели пептидного остатка
из формулы каждой аминокислоты вычитается одна молекула
воды, то есть $10$ протонов и $8$ нейтронов.
Цвиттер-ионная форма имеет тот же атомный состав,
что и свободная нейтральная молекула, поэтому её нуклонные
числа совпадают с исходными.

Эти изменения описываются однородным целочисленным сдвигом
$(s_P,s_N)$: к числам каждой аминокислоты прибавляется
одна и та же пара значений. В разности двух сумм одной
величины добавки сокращаются, если число слагаемых
одинаково. Поэтому сохраняются разности молекулярных
чисел триптофана и изолейцина,
$\delta N(\mathrm{Trp},\mathrm{Ile})=37$ и
$\delta P(\mathrm{Trp},\mathrm{Ile})=36$.
Сохраняются и разности сумм Кето/Амино
и Сильные/Слабые третьей позиции: в пуле без служебных
кодонов каждая из этих групп содержит по~$30$ кодонов.
Это устойчивость относительно однородных сдвигов;
произвольное изменение молекулярного представления
может менять разные аминокислоты по-разному.

В абсолютных суммах одинаковые добавки не сокращаются.
Делимость нейтронной суммы пула на~$37$ сохраняется
при условии $60s_N\equiv0\pmod{37}$, то есть
$s_N\equiv0\pmod{37}$, поскольку $60$ и $37$
взаимно просты. Для пептидных остатков нейтронная
и протонная суммы равны соответственно
\[
\begin{aligned}
3589-60\cdot8&=3109=84\cdot37+1,\\
4180-60\cdot10&=3580=97\cdot37-9.
\end{aligned}
\]
Обе теряют исходные остатки по модулю~$37$.
Нейтронные суммы Кето и Амино третьей позиции становятся
равными $1573$ и $1536$: разность~$37$ сохраняется,
но ни одна сумма больше не кратна~$37$.
Нуклонная сумма Октета~I в том же представлении
равна $3700-32\cdot18=3124$, поэтому её разность
с нейтронной суммой пула составляет $3124-3109=15$
вместо исходных~$111$.

Таким образом, абсолютные арифметические свойства
относятся к принятому молекулярному представлению.
Разности между равномощными группами пула и между молекулами
сохраняются при проверенных однородных сдвигах,
но из этого не следует устойчивость всей структуры.
Все $13$ проверяемых величин при трёх представлениях
приведены в \texttt{representation\_sensitivity.csv}.

\paragraph{Асимметрия Кето/Амино по трём позициям.}
Сравним теперь разности сумм Кето минус Амино
в пуле без служебных кодонов. В стандартном коде
на каждой позиции ровно одна из четырёх разностей
по $T$, $P$, $N$ и $\Delta$ кратна~$37$:
\[
\begin{array}{c@{\qquad}ccc}
\text{Позиция} & 1 & 2 & 3 \\
\text{Величина} & N & T & N \\
\text{Разность} & -333=-9\cdot37 & 185=5\cdot37 & 37
\end{array}
\]
На второй позиции делимость возникает при вычитании:
сами нуклонные суммы не кратны~$37$, но имеют одинаковые
остатки, которые сокращаются в разности.

На первых двух позициях в разностях остаются вклады
всех девятнадцати аминокислот пула (метионин выпадает вместе с $\mathrm{ATG}$). На третьей благодаря
синонимии сокращаются вклады всех аминокислот, кроме
триптофана и изолейцина. Нейтронная разность двух групп
тем самым сводится к разности одной пары молекул
(раздел~\ref{sec:results-key1-trpile}):
\[
N_{\mathrm{SE}}(\mathrm{Keto})
-N_{\mathrm{SE}}(\mathrm{Amino})
=N(\mathrm{Trp})-N(\mathrm{Ile})=37.
\]

Сравнение $27$ таблиц NCBI показывает, что сочетание
делимостей по $N$, $T$, $N$ на трёх позициях
воспроизводится только у таблиц~$1$, $11$ и $28$
в Модели~1 и у таблиц~$1$ и $11$ в Модели~2.
Во всех этих случаях состав анализируемого пула
и его аминокислотные назначения совпадают со стандартными.
У таблицы~$28$ такое совпадение получается при исключении
контекстно-зависимых стоп-кодонов в Модели~1;
Модель~2 возвращает их в пул
(раздел~\ref{sec:discussion-trpile-robustness}).
В остальных таблицах делимость для этой последовательности
сохраняется не более чем на одной позиции.
Полные результаты приведены в
\texttt{position\_asymmetry\_ncbi.csv}.

Последовательность $N$, $T$, $N$ выделена после просмотра данных.
Сравнение таблиц описывает распространённость этого сочетания
в реестре и не даёт оценки вероятности его случайного возникновения.
Все четыре набора контрольных данных вычисляются
детерминированно в точной целочисленной арифметике
скриптом \texttt{controls.py}
(раздел~\ref{sec:datasets}).

\subsection{Значимость при нулевой модели случайного кода}
\label{sec:discussion-montecarlo}

Сравнение с природными кодами и детерминированные контроли
(разделы~\ref{sec:discussion-trpile-robustness}
и~\ref{sec:discussion-controls}) описывают специфичность
наблюдаемой структуры. Сравнение природных кодов само по себе
не определяет вероятность наблюдаемых свойств при случайных
перестановках аминокислот. Здесь оценивается, как часто при случайном
перераспределении аминокислот между блоками стандартного кода
возникает сопоставимая концентрация групповых величин на
кратных~$37$ и вблизи них. Такое сравнение со случайными альтернативами применяется
в исследованиях оптимальности генетического кода~\citep{freeland1998}.

Нулевая модель сохраняет $64$ кодона, двадцать блоков синонимов
с их размерами, три стоп-кодона, разбиение Румера и химические
оси третьей позиции. В каждом испытании двадцать реальных
аминокислот равномерно случайно переставляются между этими
блоками; протонное и нейтронное числа каждой аминокислоты
переносятся вместе. Рассматриваются шесть основных сводок.
\texttt{S1}--\texttt{S3} и \texttt{S4sense} вычисляются в пуле
без служебных кодонов ($60$ кодонов): из каждой группы в каждом испытании
исключаются $\mathrm{ATG}$, $\mathrm{TAA}$, $\mathrm{TAG}$ и
$\mathrm{TGA}$. \texttt{S4a} и \texttt{S4b} распространяют
анализ на полную таблицу при Key~0 и Key~1 соответственно.

Для каждой из шести основных сводок по $M=10^{6}$ испытаниям получена оценка
хвостовой вероятности $p_{\mathrm{MC}}=(b+1)/(M+1)$, где $b$~---
число случайных кодов с результатом не менее экстремальным, чем
наблюдаемый. Для счётчиков делимости учитываются равные или
большие значения, для среднего расстояния~--- равные или меньшие.
Вместе с оценкой приводится $95\%$-й интервал Уилсона для доли
$b/M$, характеризующий неопределённость Монте-Карло.
Расчёт реализован в \texttt{mc\_significance.py} с начальным
значением генератора \texttt{seed=0}; обозначения
сводок и результаты приведены в \texttt{mc\_significance.csv}.
Независимая реализация \texttt{verify\_mc.py} использует другое
построение кодонных групп, прямой перебор при выводе Key~1 и
другое начальное значение генератора (\texttt{seed=12345}).
Она воспроизводит наблюдаемые
статистики точно, а хвостовые вероятности~--- в пределах ожидаемой
ошибки Монте-Карло.

\texttt{S1} проверяет совместную делимость на~$37$ нейтронной
суммы пула без служебных кодонов и нуклонной суммы Октета~I.
В стандартном коде это
$N(\mathcal{C}_{\mathrm{SE}}) = 3589 = 97 \cdot 37$ и
$T(\mathrm{Octet~I}) = 3700 = 100 \cdot 37$.
В случайных кодах проверяется делимость обеих сумм, без требования
совпадения с их точными значениями в стандартном коде.
Получено $p_{\mathrm{MC}} = 0{,}000800$ с $95\%$-м интервалом
$0{,}000746$--$0{,}000856$. Совместное событие оценивается
непосредственно, без предположения о независимости двух делимостей.

\texttt{S2} подсчитывает делящиеся на~$37$ величины $T$, $P$,
$N$, $\Delta$ для девяти групп, ограниченных пулом без служебных
кодонов: полного пула, шести половин химических осей и двух
октетов. В стандартном коде делятся $13$ из $36$ величин,
тогда как среднее в случайных кодах равно $1{,}0299$.
Результат не меньше $13$ получен в $39$ из $10^{6}$ испытаний:
$p_{\mathrm{MC}} = 4{,}0 \times 10^{-5}$,
$95\%$-й интервал $2{,}9 \times 10^{-5}$--$5{,}3 \times 10^{-5}$.
Величины связаны тождествами $T=P+N$, $\Delta=P-N$ и
попарным равенством сумм групп-близнецов Кето/Сильные и
Амино/Слабые; суммы взаимодополняющих групп дают сумму всего пула.
Эти зависимости сохраняются в каждой перестановке
и уже учтены при сравнении со случайными кодами.

С учётом этих тождеств наблюдаемые $13$ делимостей сводятся
к кратности~$37$ пяти величин: $N(\mathcal{C}_{\mathrm{SE}})$,
$N_{\mathrm{SE}}(\mathrm{Keto})$, $P_{\mathrm{SE}}(\mathrm{Amino})$,
$\Delta(\mathrm{Pyrimidine})$ и $T(\mathrm{Octet~I})$.
Это пять линейно независимых сравнений по модулю~$37$;
статистическая независимость событий не предполагается.
Тест \texttt{S2} учитывает любые наборы из $13$ и более делимостей,
включая отличающиеся от наблюдаемого.

\texttt{S3} применяет тот же счёт только к $17$ группам,
которые изначально не содержат служебных кодонов, всего к $68$
величинам. В стандартном коде делятся $14$ величин при среднем
$2{,}1311$ в случайных кодах:
$p_{\mathrm{MC}} = 0{,}025135$,
$95\%$-й интервал $0{,}024829$--$0{,}025443$.
Одиннадцать делимостей следуют из
$T(\mathrm{Octet~I}) = 100 \cdot 37$, а остальные три образуют
пиримидиновый $\Delta$-каскад:
$\Delta(\mathrm{Pyrimidine}) = 8 \cdot 37$ и
$\Delta(\{\mathrm{C}\}) = \Delta(\{\mathrm{T}\}) = 4 \cdot 37$.
Таким образом, эта вторичная сводка отражает распространение
двух структурных каскадов.

\texttt{S4sense} оценивает общую близость групповых величин
к кратным~$37$. Для каждой из $33$ групп, ограниченных пулом
без служебных кодонов, вычисляются расстояния $T$, $P$, $N$,
$\Delta$ до ближайшего кратного~$37$; затем усредняются все
$132$ расстояния. Нулевые значения включаются с расстоянием
ноль. В стандартном коде среднее расстояние равно $4{,}8182$,
в случайных кодах~--- $9{,}1495$:
$p_{\mathrm{MC}} = 0{,}000274$,
$95\%$-й интервал $0{,}000242$--$0{,}000307$.
Тем самым нетипичность пула проявляется и в близости к кратным,
включая значения, которые не делятся на~$37$ точно.

В дополнительной апостериорной проверке (\texttt{S4pos})
из набора расстояний \texttt{S4sense} исключались нулевые
расстояния в стандартном и каждом случайном коде. Если все
расстояния нулевые, показатель по определению равен нулю;
такие испытания не пропускаются. Среднее оставшихся расстояний
равно $6{,}1154$ в стандартном коде и $9{,}4195$ в случайных:
номинальное $p_{\mathrm{MC}} = 0{,}001602$, $95\%$-й интервал
$0{,}001525$--$0{,}001681$.
Таким образом, отклонение от случайных кодов сохраняется
и после исключения точных попаданий в кратные~$37$.

В основной модели вместе с $\mathrm{ATG}$ из пула исключается
та аминокислота, которая попала в его блок. Для проверки роли
этого выбора в дополнительной условной модели метионин закреплён
за $\mathrm{ATG}$, а остальные $19$ аминокислот переставляются
между оставшимися блоками. Получено
$p_{\mathrm{MC}} = 0{,}000026$ для \texttt{S2},
$0{,}000268$ для \texttt{S4sense} и $0{,}001531$ для
\texttt{S4pos}. Нетипичность сохраняется и при фиксированном
наборе аминокислот, представленных в пуле.

Дополнительные расчёты выполнены по $10^{6}$ испытаний на модель
скриптом \texttt{mc\_sensitivity.py} (\texttt{seed=0});
полные оценки и интервалы приведены в \texttt{mc\_sensitivity.csv}.
Независимая реализация \texttt{verify\_mc\_sensitivity.py}
(\texttt{seed=12345}) воспроизводит наблюдаемые значения точно,
а средние и хвостовые вероятности~--- в пределах ожидаемой
ошибки Монте-Карло.

\texttt{S4a} вычисляет то же среднее расстояние при Key~0.
Стоп-кодоны дают нулевой вклад, а $\mathrm{ATG}$ учитывается
как аминокислота, назначенная его блоку в данном коде:
метионин в стандартном коде и результат перестановки в случайном.
Среднее расстояние составляет $5{,}5606$ в стандартном коде
и $9{,}0157$ в случайных кодах:
$p_{\mathrm{MC}} = 0{,}002807$,
$95\%$-й интервал $0{,}002704$--$0{,}002912$.
Следовательно, близость к кратным сохраняется и при таком
формальном возвращении служебных кодонов.

\texttt{S4b} оценивает среднее расстояние после вывода Key~1.
Для стандартного и каждого случайного кода заново выводится
собственное минимальное неотрицательное назначение по условиям
баланса Кето/Амино и делимости глобальных протонной и нейтронной
сумм на~$37$ (раздел~\ref{sec:results-key1-uniqueness}).
Таким образом, делимости, обеспечиваемые самой процедурой,
возникают и в случайных кодах. Среднее расстояние равно
$4{,}6364$ в стандартном коде и $7{,}7397$ в случайных кодах:
$p_{\mathrm{MC}} = 0{,}021106$,
$95\%$-й интервал $0{,}020825$--$0{,}021389$.
Эта вероятность характеризует близость групповых величин
к кратным~$37$ после одинаковой процедуры вывода ключа
для всех сравниваемых кодов.

При выбранной модели случайного перераспределения аминокислот
стандартный код выделяется прежде всего свойствами пула без
служебных кодонов. В нём больше групповых величин точно делятся
на~$37$, а среднее расстояние до ближайшего кратного меньше,
чем обычно в случайных кодах. После исключения точных попаданий
оставшиеся величины также расположены в среднем ближе к кратным.
Следовательно, наблюдаемая близость не исчерпывается самими
делимостями. Число делимостей и общее среднее расстояние связаны:
точные делимости входят в это среднее с нулевым расстоянием.
Эти результаты характеризуют разные стороны одной арифметической
структуры. Закрепление метионина за $\mathrm{ATG}$ также
не устраняет эту нетипичность. При возвращении служебных кодонов в рамках двух
формальных назначений близость к кратным сохраняется; после
повторного вывода Key~1 отклонение от случайных кодов выражено слабее.

Все приведённые $p_{\mathrm{MC}}$ являются номинальными
вероятностями для отдельных сводок: поправка на множественное
тестирование не применялась ни к основным сводкам, ни к дополнительным
проверкам, а зависимые
$p$-значения не объединяются перемножением. Модуль~$37$ взят
из предшествующей литературы; основные сводки формально определены до
окончательного запуска Монте-Карло, но после изучения наблюдаемой
структуры. Анализ не был предрегистрирован и имеет
исследовательский характер. Полученные вероятности относятся
к рассмотренным моделям с фиксированными блоками стандартного кода и
реальными нуклонными профилями аминокислот. Они не оценивают
вероятность возникновения природного кода и не указывают
на биологический механизм.

\subsection{Связь с предшествующими работами}
\label{sec:discussion-prior-work}

Структурной основой настоящего анализа служит разбиение
Румера~\cite{rumer1966,rumer1968} на восемь полностью вырожденных
и восемь остальных кодонных семейств, образующих два октета.
Арифметические свойства нуклонных сумм, в том числе их делимость
на~$37$, подробно рассматриваются у Щербака и
Макукова~\cite{makukov2013} и у Панова и Филатова~\cite{panov2024}.
В обеих работах используются разные способы подсчёта: по отдельным
кодонам и по сокращённым кодонным сериям. Поэтому связь с их
результатами требует явного указания, какие вклады суммируются.

При сокращённом подсчёте каждый аминокислотный продукт учитывается
один раз в каждом семействе, независимо от числа кодирующих его
кодонов внутри этого семейства; повторы между семействами
сохраняются. Для Октета~I сумма полных нуклонных чисел при таком
подсчёте равна $925=25\cdot37$. Для Октета~II при нулевом вкладе
стоп-кодонов она равна $2220=60\cdot37$; в разложении на постоянные
и боковые части это $1110+1110$
\cite{makukov2013,panov2024}. Эти суммы относятся к сокращённому
представлению, тогда как обозначение $T$ групп в настоящей работе
всегда предполагает подсчёт по кодонам.

При подсчёте по кодонам Панов и Филатов~\cite{panov2024}
получают для двух октетов суммы $3700$ и $4218$. Стоп-кодонам
в этом расчёте приписан ноль, а $\mathrm{ATG}$ учитывается как
метионин, что соответствует нашему Key~0:
\[
T(\mathrm{Octet~I})=3700=100\cdot37,
\qquad
T_{0}(\mathrm{Octet~II})=4218=114\cdot37.
\]
Предложенное ими альтернативное назначение каждому стоп-кодону
полного нуклонного числа~$74$ увеличивает вторую сумму до
$4218+3\cdot74=4440=120\cdot37$; первая остаётся неизменной.

Для Октета~I связь двух способов подсчёта задаётся четырёхкратной
вырожденностью: $3700=4\cdot925$. Каждая из его $16$-кодоновых
половин химических осей имеет сумму $1850$, а каждая
$8$-кодоновая группа с фиксированным третьим основанием~--- $925$.
В Октете~II при переходе к кодонной сумме каждый вклад требуется
учитывать с кратностью соответствующей кодирующей серии.
Однако в стандартном коде сокращённая сумма совпадает с двумя
суммами при Key~0:
\[
T_{0}(\mathrm{Keto}\cap\mathrm{Octet~II})
=T_{0}(\mathrm{Strong}\cap\mathrm{Octet~II})=2220.
\]
В каждой из этих половин каждый аминокислотный продукт своего
семейства представлен одним кодоном, а стоп-кодоны дают нулевой
вклад. Таким образом, совпадение определяется фактическими
назначениями кодонов внутри семейств. Распространение этих
делимостей на другие группы при Key~0 разобрано в
разделе~\ref{sec:results-key0}.

Разложение молекулы на постоянную часть и боковую цепь играет
существенную роль в обеих предшествующих работах. У Щербака и
Макукова операция \emph{activation key}~\cite{makukov2013}
формально переносит один
нуклон из боковой цепи пролина в его постоянную часть: числа
$42$ и $73$ заменяются на $41$ и $74$. Полное нуклонное число
молекулы при этом сохраняется: $42+73=41+74=115$.
В настоящей работе используются полные свободные нейтральные
молекулы; их протонные и нейтронные числа не требуют такого
перераспределения между частями молекулы.

Наш анализ рассматривает $33$ фиксированные группы, построенные
на октетах Румера и химических осях третьей позиции, и для каждой
группы последовательно сопоставляет $T$, $P$, $N$ и
$\Delta=P-N$. В пуле без служебных кодонов нейтронная разность
$37$ и протонная разность $36$ в паре $\mathrm{Trp}$--$\mathrm{Ile}$
характеризуют асимметрию Кето/Амино
(раздел~\ref{sec:results-localization}). Их раздельный учёт
вместе с глобальными суммами лежит в основе вывода служебного
назначения из условий баланса и делимости
(раздел~\ref{sec:results-key1-derivation}).
Совпадение полного вклада каждого стоп-кодона с предложенным
ранее значением~$74$ не делает это назначение тождественным
Key~1: в расчёте Панова и Филатова сохраняется метиониновый вклад
$\mathrm{ATG}$, равный $149$, тогда как при Key~1
$T_{1}(\mathrm{ATG})=1$. При этом наш ключ задаёт отдельно
протонный и нейтронный вклады каждого служебного кодона.

\subsection{Ограничения и перспективы}
\label{sec:discussion-limitations}

Статистическая значимость наблюдаемых арифметических
закономерностей оценивается в
разделе~\ref{sec:discussion-montecarlo} относительно основной
нулевой модели, которая фиксирует архитектуру вырожденности кода
и нуклонные числа аминокислот и случайно переставляет аминокислоты
между кодонными семействами. Дополнительная условная модель
закрепляет метионин за $\mathrm{ATG}$, сохраняя набор аминокислот,
представленных в пуле без служебных кодонов. Наблюдаемые счёт
делимостей (\texttt{S2}) и среднее расстояние до ближайшего
кратного~$37$ (\texttt{S4sense}) нетипичны при обеих моделях.
Эти сравнения имеют исследовательский характер и зависят от
выбора групп и показателей, в том числе от повторного учёта
связанных величин. Они не оценивают вероятность обнаружить
произвольную арифметическую закономерность после поиска по данным.
Более узкий тест \texttt{S3}, ограниченный вложенными группами,
которые вообще не содержат служебных кодонов, даёт лишь слабую
нетипичность ($p_{\mathrm{MC}} \approx 0{,}025$) и в основном
отражает копии делимости полного $T(\mathrm{Octet~I})$.
Выбор нулевой модели не единственен: другие порождающие модели,
например рандомизация самих размеров блоков или нуклонных чисел,
исследовали бы иные аспекты этой закономерности и оставлены для
будущей работы, как и любое объяснение механизма, способного её
породить.

Параметризация служебных кодонов в настоящей
работе ограничена Key~0 и Key~1;
дополнительные параметрические назначения
приведены в \texttt{key\_parameters.csv} и в
сопровождающем \texttt{visualization.html}
для иллюстрации других вариантов выбора.

Дальнейшее направление связано с группо-теоретической
структурой таблицы кодонов. Следуя~\citet{panov2024},
каждому из $24$ упорядочений четырёх оснований на
осях кодонной матрицы сопоставляется его
\emph{каллиграмма}: геометрический паттерн,
образуемый восемью однородными кодонными семействами
Октета~I на матрице
$4 \times 4$. Связность октета относится к
связности этого паттерна относительно
$4$-смежности клеток матрицы. Среди $24$
каллиграмм Панoв и Филатов~\cite{panov2024}
выделили четыре ---
$\mathrm{CTGA}$, $\mathrm{CGTA}$, $\mathrm{ATGC}$
и $\mathrm{AGTC}$ --- в которых каждый из
румеровских октетов образует связную область, и показали, что
эти четыре образуют замкнутое множество относительно
преобразований
$R = (\mathrm{T} \leftrightarrow \mathrm{G},
\mathrm{C} \leftrightarrow \mathrm{A})$,
$R_1 = (\mathrm{T} \leftrightarrow \mathrm{G})$, и
$R_2 = (\mathrm{C} \leftrightarrow \mathrm{A})$.
Эти три преобразования вместе с тождественным
преобразованием образуют группу Клейна $V_4$, а
действие $V_4$ естественным образом распространяется
на полное множество из $24$ каллиграмм. Это действие
$V_4$ на полном множестве из $24$ каллиграмм является
свободным, разбивая его на шесть орбит по четыре
каллиграммы в каждой. Ровно одна из этих шести орбит
состоит из связных каллиграмм, что и даёт четыре
указанные выше каллиграммы; остальные пять орбит не
содержат ни одной связной каллиграммы. Более того,
эта связная орбита содержит канонический химический
порядок $C, G, T, A$ Румера~\cite{rumer1968}
(раздел~\ref{sec:genetic-code}); соответствие
между этим химическим порядком и геометрическим
свойством связности не является автоматическим,
поскольку лишь два из восьми химически чередующихся порядков четырёх оснований дают связные октеты. Поиск дальнейших
структурных инвариантов, различающих орбиты, а также
связей между группо-теоретическими орбитами и
арифметическими закономерностями настоящей работы,
оставляется для будущих исследований.

\section{Заключение}
\label{sec:conclusion}

Стандартный генетический код демонстрирует арифметические регулярности, проявляющиеся на нескольких уровнях параметрической специфичности. Первый уровень не зависит ни от какого выбора нуклонных вкладов для служебных кодонов и содержит центральный результат данной работы: при суммировании по шестидесяти кодонам $\mathcal{C}_{\mathrm{SE}}$ нейтронный пул оказывается точным кратным модулю, $N(\mathcal{C}_{\mathrm{SE}}) = 3589 = 97 \cdot 37$, тогда как протонный пул не достигает на одну единицу нечётного кратного, $P(\mathcal{C}_{\mathrm{SE}}) \equiv -1 \pmod{37}$---свойство протонно-нейтронного разделения, которое не проявляется на уровне полного числа $T$. Наряду с этим существуют ещё два факта, не зависящих от параметризации. Восемь семейств Октета~I содержат $T(\mathrm{Octet~I}) = 3700 = 100 \cdot 37$, при $P(\mathrm{Octet~I}) = 2000$ и $N(\mathrm{Octet~I}) = 1700$, причём последняя величина вместе с $N_{\mathrm{SE}}(\mathrm{Octet~II})$ образует сумму $1700 + 1889 = 3589$; кроме того, вся асимметрия Кето/Амино в коде после исключения служебных кодонов локализуется в единственной паре аминокислот $\mathrm{Trp}$--$\mathrm{Ile}$, для которой $|\delta N| = 37$ и $|\delta P| = 36$. Универсальная пиримидиновая синонимия внутри Октета~II порождает тождества между двумя шестнадцатикодонными группами $\mathrm{Keto}$/$\mathrm{Strong}$ и $\mathrm{Amino}$/$\mathrm{Weak}$ внутри Октета~II.

Второй слой возникает при наивном прочтении, когда числа протонов и нейтронов
$P$ и $N$ анализируются отдельно. Суммарное число нуклонов
$T$ делится на~$37$ для четырех
групп химических осей и для Октета~II.
Для групп $\mathrm{Amino}$ и $\mathrm{Weak}$
значения $P$ и $N$ делятся на~$37$ по отдельности,
что выравнивает эти группы с решеткой шага~$37$
вдоль обеих осей.

Третий уровень включает назначение служебных кодонов. Существует единственное минимальное неотрицательное целочисленное назначение, Key~1 (раздел~\ref{sec:results-key1-derivation}), при котором параметризационно-независимая структура распространяется на полную таблицу из $64$ кодонов. Key~1 не является основанием для приведённых выше результатов, а представляет собой средство, позволяющее распространить их на полную таблицу: та же делимость для пула без служебных кодонов, та же локализованная пара и та же протонная чётность, которые выполняются при неназначенных служебных кодонах, именно и делают возможным нахождение решения в малых целых числах. Нейтронный баланс вдоль оси Кето/Амино воспроизводит разность $\mathrm{Trp}$--$\mathrm{Ile}$, $|\delta N| = 37$; глобальный нейтронный дефицит $T(\mathrm{Octet~I}) - N(\mathcal{C}_{\mathrm{SE}}) = 111 = 3 \cdot 37$ замыкается, поскольку является общим кратным числа стоп-кодонов и модуля; теорема о чётности числа протонов исключает возможность $|\delta P| = 37$ и не допускает, чтобы протонная сумма пула без служебных кодонов оказалась на нечётных кратных модуля, вследствие чего она располагается на одну единицу ниже них---причём благодаря конкретным нуклонным значениям этот отступ составляет в точности одну единицу---и этот дефицит компенсируется минимальным вкладом старт-кодона $p_{\text{start}} = 1$; а условие $\gcd(4, 37) = 1$ оставляет единственное минимальное решение. Ни одно из этих свойств не является параметрическим предположением; каждое представляет собой эмпирическое свойство стандартного генетического кода, установленное на параметризационно-независимом уровне, а Key~1 существует как минимальное расширение именно потому, что все они выполняются одновременно.

При выведенной параметризации арифметические закономерности
организуются в связное семейство: выравнивание четырех
групп химических осей вдоль осей $T$, $P$, $N$ и $\Delta$
(разделе~\ref{sec:results-key1-axisalignment}),
макроскопические тождества, включая
$N(\mathrm{All}) = T(\mathrm{Octet~I}) = 3700$
и $P(\mathrm{All}) = T(\mathrm{Octet~II}) =
4292$
(разделе~\ref{sec:results-key1-macroidentities}),
единый инвариант $\Delta = 296$ для всех шести
групп из $32$ кодонов
(разделе~\ref{sec:results-key1-deltainvariant}), и
паттерны делимости на~$37$ и~$999$ на множественных
уровнях разбиения
(разделе~\ref{sec:results-key1-divisibility37and999}).
Сравнительный анализ по $27$ таблицам трансляции
NCBI (разделе~\ref{sec:results-key1-ncbi}) показывает,
что одновременное выполнение всех перечисленных выше
структурных предпосылок является свойством исключительно
стандартного кода.

Две особенности этой структуры ограничивают то, как
её можно трактовать. Во-первых, закономерности
заключены в групповых суммах, а не в отдельных
молекулах: ни одна отдельная каноническая
аминокислота не имеет нуклонного количества,
делящегося на~$37$, а в группах с переменной
вырожденностью кодирования, таких как
$\mathrm{Amino}$, делимость выполняется только при
счёте кодонов с их кратностью и нарушается, если
каждая аминокислота учитывается однократно. Таким
образом, нуклонный состав и вырожденность
кодирования не являются в этой структуре
независимыми величинами: закономерность есть
совместное свойство того, как распределена нуклонная
масса и сколько кодонов получает каждая
аминокислота. Во-вторых, весь первый уровень,
изложенный выше, включая значения по
Октету~I и тождества между
группами-близнецами, устанавливается при служебных
кодонах, не несущих нуклонных значений вовсе;
выведенная параметризация служебных кодонов входит
лишь на последующих уровнях и не влияет на
параметр-независимые результаты. Приписанные значения
суть, следовательно, приём для предъявления
арифметики полной таблицы, а не основание, на котором
держатся основные результаты.

Специфичность, установленная сравнением с природными
кодами, дополняется детерминированными контролями и
анализом Монте-Карло. Контроли показывают, что делимости
сосредоточены в группировках по третьей позиции и что
$37$ выделяется среди проверенных простых модулей.
При явной нулевой модели, фиксирующей архитектуру вырожденности
и переставляющей двадцать реальных протонно-нейтронных профилей
по двадцати блокам, две взаимозависимые сводки пула без
служебных кодонов помещают стандартный код в дальний хвост:
дискретный счёт \texttt{S2} даёт
$p_{\mathrm{MC}} = 4{,}0 \times 10^{-5}$, а среднее расстояние
\texttt{S4sense} --- $p_{\mathrm{MC}} = 2{,}74 \times 10^{-4}$.
Более узкий тест \texttt{S3}, ограниченный группами, которые
вообще не содержат служебных кодонов, даёт слабую нетипичность
($p_{\mathrm{MC}} \approx 0{,}025$) и в основном отражает зависимые
копии делимости $T(\mathrm{Octet~I})$. При Key~0 концентрация
сохраняется ($p_{\mathrm{MC}} \approx 0{,}0028$), тогда как
симметричная проверка процедуры вывода Key~1 даёт умеренную
нетипичность ($p_{\mathrm{MC}} \approx 0{,}021$). Эти вероятности
относятся к разным, зависимым сводкам и не объединяются
перемножением.
Дополнительные апостериорные проверки показывают, что отклонение
сохраняется после исключения нулевых расстояний и при закреплении
метионина за $\mathrm{ATG}$. Эти результаты характеризуют
нетипичность в рассмотренных комбинаторных моделях, отличную от
специфичности среди $27$ природных кодов, и не указывают на
биологический механизм. Другие нулевые модели оставлены для
будущей работы.

Методологический выбор, лежащий в основе настоящей работы,
заключается в том, чтобы ограничить внимание параметрами,
определение которых однозначно — суммарным числом
нуклонов~$T$, его разложением на протоны~$P$ и
нейтроны~$N$ и разностью $\Delta = P - N$ —
и единой системой группировки, основанной на октетах
Румера и химических осях третьей позиции.
Интерпретация собранных здесь арифметических фактов
остается открытой.

\bibliography{references}

\end{document}
